\documentclass[11pt]{amsart}

\usepackage[margin=1.2in]{geometry}
\usepackage{amsmath,amssymb,amsthm}
\usepackage{graphicx}
\usepackage{xcolor}
\definecolor{linkblue}{rgb}{0.1,0.1,0.5}
\usepackage[colorlinks=true,linkcolor=linkblue,citecolor=linkblue,urlcolor=linkblue]{hyperref}

\theoremstyle{plain}
\newtheorem{theorem}{Theorem}[section]
\newtheorem{proposition}[theorem]{Proposition}
\newtheorem{lemma}[theorem]{Lemma}
\newtheorem{corollary}[theorem]{Corollary}
\newtheorem{conjecture}[theorem]{Conjecture}
\newtheorem{mainthm}{Theorem}
\renewcommand{\themainthm}{\Alph{mainthm}}
\newtheorem*{theoremA}{Theorem A}

\theoremstyle{definition}
\newtheorem{definition}[theorem]{Definition}
\newtheorem{example}[theorem]{Example}
\newtheorem{problem}[theorem]{Problem}

\theoremstyle{remark}
\newtheorem{remark}[theorem]{Remark}

\newcommand{\HH}{\mathbb{H}}
\newcommand{\CC}{\mathbb{C}}
\newcommand{\RR}{\mathbb{R}}
\newcommand{\ZZ}{\mathbb{Z}}
\newcommand{\NN}{\mathbb{N}}
\newcommand{\CP}{\mathbb{CP}}
\newcommand{\HP}{\mathbb{HP}}
\newcommand{\RP}{\mathbb{RP}}
\renewcommand{\SS}{\mathbb{S}}
\newcommand{\Xii}{\Xi}
\newcommand{\re}{\operatorname{Re}}
\newcommand{\im}{\operatorname{Im}}
\newcommand{\id}{\operatorname{id}}
\newcommand{\Aut}{\operatorname{Aut}}
\newcommand{\LP}{\mathcal{LP}}
\newcommand{\EE}{\mathcal{E}}
\newcommand{\dd}{\,d}

\begin{document}

\title[The zero-sphere program]{The zero-sphere program: a quaternionic--twistor framework for the Riemann Hypothesis}

\author{Edward B. Baker III}
\email{edwardbaker86@gmail.com}

\date{\today}

\begin{abstract}
We describe a research program on the Riemann Hypothesis organized around the canonical extension of
the completed zeta function $\Xi$ to the quaternions as an intrinsic slice regular function. The zero
set of this extension consists of real points and round $2$-spheres, and RH is equivalent to the
absence of spheres. We develop three faithful formulations of this dichotomy: a twistor-geometric one
(the zero divisor of an induced self-map of an incidence quadric is invariant under the real structure
component-by-component), a positivity hierarchy graded by spherical Bessel tags whose bottom rung is
the Laguerre inequality, and an explicit-formula criterion --- a numerically certified,
prime-computable quadratic form whose boundedness in every frequency window is an unconditional
consequence of RH, and which, under a mild non-accumulation hypothesis on the zeros, grows at
exponential rate exactly twice the supremal sphere radius, making RH equivalent to
subexponential growth. Two no-go theorems delimit the
framework: no continuous linear functional of a rotational extension, and no deformation-invariant
degree count, distinguishes real zeros from spheres. The organizing statement is that RH asserts the
agreement, on the zero set, of the holomorphic involution $\rho\mapsto 1-\rho$ with the
antiholomorphic involution $\rho\mapsto\bar\rho$; the central open problem is to import the latter
into explicit-formula arithmetic. The Davenport--Heilbronn function serves as control throughout, and
all numerical claims are reproducible from the included code. This is a program paper: it proves
structure theorems and delimiting results, states one sharp organizing question, and claims no
progress on RH itself.
\end{abstract}

\maketitle
\tableofcontents

%=====================================================================
\section{Introduction}\label{sec:intro}

Rotate the critical line onto the real axis. Writing $s=\tfrac12+iz$ and
\[
\xi(s)=\tfrac12 s(s-1)\pi^{-s/2}\Gamma(s/2)\zeta(s),\qquad
\Xi(z):=\xi\!\left(\tfrac12+iz\right),
\]
the function $\Xi$ is entire, even, of order $1$, and real on $\RR$, and the Riemann Hypothesis is
the statement that all zeros of $\Xi$ are real. Because the Taylor coefficients of $\Xi$ at the
origin are real, $\Xi$ extends canonically to the quaternions $\HH$ as an \emph{intrinsic slice
regular} function, and by the zero-structure theorem for such functions its zero set is a union of
isolated real points and isolated round $2$-spheres $x+y\,\SS$, where $\SS$ is the sphere of
imaginary units. A real zero of the complex $\Xi$ gives a point; a conjugate pair $a\pm iy_0$ of
nonreal zeros gives the sphere $a+y_0\SS$. Hence:

\begin{theoremA}[Reformulation; proved as Theorem~\ref{thm:reformulation} in \S\ref{sec:prelim}]
RH holds if and only if the zero set of the quaternionic extension
$\Xi\colon\HH\to\HH$ contains no $2$-sphere.
\end{theoremA}

By itself this is a relabeling, not progress, and we say so plainly: every statement in this paper
should be read against the standing question \emph{``what does the reformulation buy?''}. The
program's answer is that the sphere/point dichotomy admits three faithful and interlocking
formulations --- geometric, positivity-theoretic, and arithmetic --- each of which exposes structure
that the complex formulation leaves implicit, and that the framework comes with built-in
falsification tests that delimit what can and cannot work. Concretely, the paper proves:

\begin{itemize}
\item two structure theorems (Theorems~\ref{thm:lift} and \ref{thm:dichotomy}): every intrinsic
slice regular function lifts to a holomorphic, equivariant self-map of an incidence quadric
$Q\cong\CP^1\times\CP^1$ inside twistor space $\CP^3$, and on $Q$ the zero divisor of the lifted
$\Xi$ distinguishes real zeros from zero spheres \emph{equivariantly} --- by the action of the real
structure on its components --- though not homologically; RH becomes the statement that every
component of the zero locus in the quotient $Q/\sigma$ is nonorientable;
\item two no-go theorems (Theorems~\ref{thm:linear-nogo} and \ref{thm:degree-nogo}): no continuous
linear functional of any rotational extension of $\Xi$, and no deformation-invariant degree count,
can detect spheres --- delimiting results that any candidate mechanism must respect;
\item one certified criterion (Theorem~\ref{thm:rate}): a prime-computable quadratic form
$|Z(c,\lambda)|^2$ built from a Gaussian test family in the explicit formula. Unconditionally, RH
implies boundedness in every frequency window, so exponential growth in any window disproves RH;
under a mild non-accumulation hypothesis on the zeros, the growth rate equals exactly twice the
supremal sphere radius and RH is equivalent to subexponential growth in every window. The explicit
formula underlying it is certified numerically to $10^{-8}$--$10^{-10}$ against $250$ zeros, and
the criterion recovers the radius of the first Davenport--Heilbronn sphere to four decimal places
from zero-side data;
\item one organizing question (Section~\ref{sec:involutions}): the zero multiset carries the
holomorphic involution $\rho\mapsto1-\rho$ (functional equation), which the explicit formula sees
natively, and the antiholomorphic involution $\rho\mapsto\bar\rho$ (reality), which defines the
spheres and which the explicit formula is structurally blind to. RH is the statement that the two
involutions agree on the zero set. The program's central problem is to import the antiholomorphic
involution into explicit-formula arithmetic --- to find unconditional identities or bounds for
\emph{conjugate-paired sums} $\sum_\gamma F(\gamma,\bar\gamma)$.
\end{itemize}

Two hardness checks run through everything. First, the Davenport--Heilbronn function
\cite{DavenportHeilbronn1936a,DavenportHeilbronn1936b,Titchmarsh1986} shares the symmetry and growth
of $\Xi$ but has zeros off the critical line --- genuine zero spheres, computed explicitly in
Section~\ref{sec:hardness} --- and no Euler product; any mechanism that would exclude spheres from
symmetry and growth alone is therefore wrong, and the arithmetic input is not optional. Second, the
nonnegativity of the de Bruijn--Newman constant \cite{RodgersTao2020} places $\Xi$ at the exact
boundary of the sphere-free regime, so any true criterion must degrade only polynomially (or worse)
as the sphere radius shrinks to zero; a soft argument stable under small perturbations proves too
much. We treat these checks as specifications, and each proposed structure in this paper is tested
against both.

\subsection*{Provenance}
The program synthesizes two threads. The first is a sequence of single-author manuscripts
(2015--2026) that searched for a quaternionic function class adapted to the rotated $\Xi$: a 2015
manuscript \cite{Baker2015} built a Fueter-regular extension hierarchy with an
$SU(2)$-representation-theoretic decomposition, raising/lowering operators, spherical Bessel
kernels, and a quaternionic Mellin inversion theory, together with a negative result on linear
pairings that Section~\ref{sec:nogo} promotes to a theorem; a companion 2016 manuscript
\cite{Baker2016} developed the associated quaternionic exponential and plane-wave expansion and
connected it to the Clifford--Fourier kernel of Clifford analysis
\cite{BrackxDeSchepperSommen2005,DeBieXu2011,BoryReyesShapiro2010}; manuscripts from 2017--2024
\cite{Baker2017,Baker2023,Baker2022arXiv,Baker2023arXiv} developed twistor-geometric machinery ---
a rigidity theorem for fiberwise degree-one self-maps of twistor space, and generalized conformal
morphisms defined on twistor correspondence spaces, applied there to Yang--Mills symmetries in the
tradition of \cite{MasonWoodhouse1996}. The
second thread is the modern theory of slice regular functions (Gentili--Struppa
\cite{GentiliStruppa2007} and its development
\cite{GentiliStoppato2008,CGSS2009,GSSbook,GhiloniPerotti2011}), whose intrinsic subclass turns out
to be exactly the function class the first thread was searching for, and whose twistor branch
\cite{GSS2014,Altavilla2018,AltavillaSarfatti2019} realizes the intended geometry. The synthesis,
the specific computations, and the numerical laboratory reported here were carried out in August
2026; the historical manuscripts are archived unaltered at \cite{ZenodoArchive}.

The role of the twistor correspondence here instantiates a pattern from the Yang--Mills thread
\cite{Baker2022arXiv,Baker2023arXiv}: transformations that are not holomorphic (or not even
well-defined pointwise) downstairs become honest holomorphic, structure-preserving maps on a twistor
correspondence space. Theorem~\ref{thm:lift} is that phenomenon in Euclidean signature for $\HH$,
with the incidence quadric as the correspondence space.

\subsection*{What this paper is not}
It is not a proof strategy with a claimed endgame. The no-go theorems and hardness checks were
derived precisely to prevent self-deception, and Section~\ref{sec:bridge} identifies exactly where
the program currently reduces to classical difficulty: at the exponential-rate level, the growth
criterion of Theorem~\ref{thm:rate} is a Gaussian-mollified form of the classical equivalence
between zero-free regions and prime-counting error terms. The candidates for genuinely new leverage
--- the graded positivity cone, the cross-term calculus, conjugate-paired identities --- are stated
as open problems in Section~\ref{sec:involutions}, each with its hardness-check annotation.

\subsection*{Acknowledgements}
The synthesis of this program, the twistor-lift and explicit-formula computations, the numerical
experiments, and substantial portions of the exposition were carried out with Claude (Anthropic) as
a research assistant in August 2026; all mathematical claims were reviewed by the author, who takes
sole responsibility for them. Code and archival materials are available at \cite{ZenodoArchive}.

%=====================================================================
\section{Preliminaries: intrinsic slice regular functions and the reformulation}\label{sec:prelim}

Let $\HH$ denote the quaternions and $\SS=\{J\in\HH: J^2=-1\}$ the $2$-sphere of imaginary units.
For $J\in\SS$ write $\CC_J=\RR+J\RR$ for the corresponding slice, a copy of $\CC$. A domain
$\Omega\subseteq\HH$ is \emph{axially symmetric} if it is invariant under all rotations about the
real axis, i.e.\ $x+yJ\in\Omega$ for one $J\in\SS$ implies $x+yJ'\in\Omega$ for all $J'\in\SS$.

\begin{definition}[\cite{GentiliStruppa2007}]
A function $f\colon\Omega\to\HH$ on an axially symmetric slice domain is \emph{slice regular} if for
every $J\in\SS$ the restriction of $f$ to $\Omega\cap\CC_J$ is holomorphic with respect to the
complex structure given by left multiplication by $J$.
\end{definition}

Slice regular functions on axially symmetric slice domains satisfy the \emph{representation
formula} of Colombo--Gentili--Sabadini--Struppa \cite{CGSS2009}: $f$ is determined by its
restriction to a single slice, via
\begin{equation}\label{eq:repformula}
f(x+yI)=\tfrac12\bigl[f(x+yJ)+f(x-yJ)\bigr]+\tfrac{IJ}{2}\bigl[f(x-yJ)-f(x+yJ)\bigr]
\qquad(I,J\in\SS).
\end{equation}
A slice regular $f$ is \emph{intrinsic} (or slice preserving) if its power-series expansion at a
real point has real coefficients; equivalently, $f(\bar q)=\overline{f(q)}$, equivalently
\begin{equation}\label{eq:intrinsic}
f(x+yJ)=\alpha(x,y)+J\beta(x,y)\qquad\text{with }F:=\alpha+i\beta\ \text{holomorphic on
}\Omega\cap\CC_i \text{ and real on }\RR.
\end{equation}
We call $F$ the \emph{stem function} of $f$ (in the sense of Ghiloni--Perotti
\cite{GhiloniPerotti2011}); the map $F\mapsto f$ is the \emph{slice extension}, denoted $R$, and it
coincides with the rotation map of \cite[eq.~(2.5)]{Baker2015}. Since $f(q)\in\CC_{e_q}$ for
intrinsic $f$ (where $q=x+ye_q$, $e_q\in\SS$, $y>0$), pointwise products of intrinsic functions stay
in the slice: the intrinsic functions on $\Omega$ form a \emph{commutative real algebra under
pointwise operations}, restriction to a slice being an algebra homomorphism onto the intrinsic
holomorphic functions; for intrinsic functions the slice product $\star$ coincides with the
pointwise product, and composition also preserves the class \cite{GSSbook}. This pointwise algebra
closure is what distinguishes the slice extension from the Fueter-type extensions of
\cite{Baker2015}, which destroy multiplicativity; the two extension hierarchies are compared in
Section~\ref{sec:sectors}.

The zero sets of slice regular functions have a rigid structure.

\begin{theorem}[Zero structure; Gentili--Stoppato \cite{GentiliStoppato2008}; see also
{\cite[Ch.~3]{GSSbook}}]\label{thm:zerostructure}
Let $f\not\equiv0$ be slice regular on an axially symmetric slice domain $\Omega$. Then the zero
set of $f$ consists of isolated points and isolated $2$-spheres of the form $x+y\SS$ with
$x,y\in\RR$, $y\neq0$.
\end{theorem}

For an intrinsic $f$ with stem function $F$, the dictionary is exact: a real zero of $F$ gives an
isolated real point zero of $f$; a conjugate pair $a\pm iy_0$ ($y_0>0$) of zeros of $F$ gives the
zero sphere $a+y_0\SS$; and nonreal zeros of $F$ occur in conjugate pairs because $F$ is real on
$\RR$.

\subsection{The extension of \texorpdfstring{$\Xi$}{Xi}}
Riemann's integral representation (see \cite[\S10.1]{Titchmarsh1986}) reads, in the rotated
variable,
\begin{equation}\label{eq:cosine}
\Xi(z)=4\int_0^\infty\Phi(u)\cos(zu)\dd u,\qquad
\Phi(u)=\sum_{n\ge1}\bigl(2\pi^2n^4e^{9u/2}-3\pi n^2e^{5u/2}\bigr)e^{-\pi n^2e^{2u}},
\end{equation}
where $\Phi$ is even, positive, and decays like $e^{-\pi e^{2u}}$; the positivity and decay are due
to P\'olya (see \cite{Polya1926,Polya1927} and the summary in
\cite[Theorem~A]{CsordasNorfolkVarga1986}). The normalizing constant $4$ in \eqref{eq:cosine}
depends on conventions for $\Phi$; with the convention displayed it is certified numerically in
Appendix~\ref{app:numerics} (the identity \eqref{eq:cosine} holds to ten digits at
$x\in\{0,5,20,30\}$, and the right-hand side vanishes to the quadrature floor at the first zero
ordinate).

Since $\cos(qu)$ is an entire intrinsic slice regular function of $q$ (its power series has real
coefficients), with
\[
\cos((x+ye_q)u)=\cos(xu)\cosh(yu)-e_q\,\sin(xu)\sinh(yu),
\]
and since $|\cos(qu)|\le e^{|\im q|\,u}\cdot O(1)$ while $\Phi$ decays super-exponentially, the
integral \eqref{eq:cosine} converges absolutely for all $q\in\HH$ and defines the \emph{entire
intrinsic slice regular extension}
\begin{equation}\label{eq:Xiq}
\Xi(q)=4\int_0^\infty\Phi(u)\cos(qu)\dd u=\sum_{n\ge0}c_{2n}q^{2n},
\end{equation}
where $c_{2n}$ are the (real) Taylor coefficients of the complex $\Xi$. This is the canonical
rotation of $\Xi$ into $\HH$: no Laplacian is applied, multiplicativity is retained, and the
extension is determined by the complex $\Xi$ through the representation formula
\eqref{eq:repformula}.

\begin{mainthm}[Reformulation]\label{thm:reformulation}
The Riemann Hypothesis holds if and only if the zero set of $\Xi\colon\HH\to\HH$ defined by
\eqref{eq:Xiq} contains no $2$-sphere; equivalently, if and only if every zero of the quaternionic
$\Xi$ is a real point.
\end{mainthm}

\begin{proof}
Immediate from Theorem~\ref{thm:zerostructure} and the dictionary above: the zeros of the stem
function $\Xi|_{\CC_i}$ are the zeros of the complex $\Xi$, whose nonreal members occur in
conjugate pairs $a\pm iy_0$ and correspond exactly to the zero spheres $a+y_0\SS$ of the extension.
RH says there are no such pairs.
\end{proof}

\begin{remark}\label{rem:relabel}
Theorem~\ref{thm:reformulation} is a relabeling; it is stated as a theorem for citability, and
because the relabeling has consequences: it places the zeros of $\zeta$ inside the zero-structure
theory of a function class with complex-analysis-grade rigidity (identity principle, open mapping,
representation formula), makes the conjugate-pair symmetry a \emph{geometric} object (a round
sphere with a center and a radius), and gives the two involutions of Section~\ref{sec:involutions}
distinct carriers. Everything after this point is an attempt to convert the relabeling into
leverage, and the reader is entitled to keep score.
\end{remark}

%=====================================================================
\section{Twistor geometry: the incidence quadric and the equivariant dichotomy}\label{sec:twistor}

This section locates the intrinsic function class inside twistor geometry and proves the two
structure theorems. Throughout we use the twistor fibration
$\pi\colon\CP^3\to\HP^1\cong S^4$, restricted over the affine chart $\HH\subset\HP^1$; our
coordinate conventions follow \cite{Baker2017,Baker2023}, and agree up to notation with those of
Gentili--Salamon--Stoppato \cite{GSS2014}. Write homogeneous coordinates
$[z_1,z_2,z_3,z_4]\in\CP^3$ with
\[
\pi[z_1,z_2,z_3,z_4]=[z_1+z_2j,\;z_3+z_4j]\in\HP^1,\qquad
q=(z_3+j)^{-1}(z_1+z_2j)\ \text{ on the chart }z_4=1 .
\]
The fiber coordinate is the stereographic parametrization of the sphere of complex structures,
\[
\eta(w)=(w+j)^{-1}\,i\,(w+j)\in\SS,\qquad w\in\CC\cup\{\infty\},
\]
and the real structure is the fixed-point-free antiholomorphic involution
\[
\sigma[z_1,z_2,z_3,z_4]=[-\bar z_2,\bar z_1,-\bar z_4,\bar z_3],
\]
whose invariant lines are exactly the fibers of $\pi$.

\subsection{Lines over spheres and the incidence quadric}

For $\zeta=x+iy\in\CC$ consider the \emph{graph line}
\[
L_\zeta=\{(x+\eta y,\ \eta):\eta\in\SS\},
\]
the tautological lift of the sphere $x+y\SS$ (for $y\neq0$), degenerating to the twistor fiber
$\pi^{-1}(x)$ when $y=0$. Substituting $\eta(w)$ and clearing denominators gives the linear
parametrization
\begin{equation}\label{eq:line}
L_\zeta=\bigl\{[\zeta w,\ \zeta,\ w,\ 1]: w\in\CP^1\bigr\},
\end{equation}
a projective line for every $\zeta\in\CP^1$, with $\sigma(L_\zeta)=L_{\bar\zeta}$; for real $\zeta$
the line is $\sigma$-invariant with the antipodal map acting on the parameter,
$w\mapsto-1/\bar w$. The Pl\"ucker image of the family is
\[
[p_{12}:p_{13}:p_{14}:p_{23}:p_{24}:p_{34}]=[\zeta^2:0:\zeta:-\zeta:0:1],
\]
a conic in the Klein quadric. The union of the family is the smooth quadric
\begin{equation}\label{eq:quadric}
Q:\ z_1z_4=z_2z_3,\qquad
Q=\{(q,\eta):q\in\CC_\eta\}\ \cong\ \CP^1_\zeta\times\CP^1_w ,
\end{equation}
the \emph{incidence variety of points with their own slice's complex structure}. Its two rulings
have transparent meaning: ruling~A ($\zeta$ constant) consists of the lines $L_\zeta$ --- lifts of
spheres, degenerating to twistor fibers over $\RR\cup\{\infty\}$; ruling~B ($w$ constant) consists
of full slices $\{(q,\eta_0):q\in\CC_{\eta_0}\}$ with a fixed complex structure. The chart identity
$q=(w+j)^{-1}\zeta(w+j)$ says that the point of $Q$ with coordinates $(\zeta,w)$ is $x+y\,\eta(w)$
with $\zeta=x+iy$: the coordinates \emph{are} (complex label of the sphere, complex structure). The
real structure restricts to
\begin{equation}\label{eq:rho}
\rho(\zeta,w)=(\bar\zeta,\ -1/\bar w),
\end{equation}
conjugation times the antipodal map, again fixed-point free.

\subsection{The lift}

\begin{mainthm}[Lift]\label{thm:lift}
Let $f$ be an intrinsic slice regular function on an axially symmetric slice domain
$\Omega\subseteq\HH$, with stem function $F$ as in \eqref{eq:intrinsic}, and let
$\Omega_Q\subseteq Q$ be the part of the quadric over $\Omega$. Then the product map
\[
\tilde f=F\times\id\colon\ \Omega_Q\longrightarrow Q,\qquad
(\zeta,w)\longmapsto(F(\zeta),\,w),
\]
is holomorphic and $\rho$-equivariant, and covers $f$: over the point $q=x+y\,\eta(w)$ of
$\Omega$ it returns $f(q)=\alpha(x,y)+\beta(x,y)\,\eta(w)$ with the same complex-structure label
$w$. If $f$ is entire, $\tilde f$ is defined on all of $Q$ minus the ruling at infinity.
\end{mainthm}

\begin{proof}
Holomorphy is immediate: $F$ is holomorphic in $\zeta$ by intrinsicness \eqref{eq:intrinsic} and
the second factor is the identity. Equivariance is the computation
$\tilde f(\rho(\zeta,w))=(F(\bar\zeta),-1/\bar w)=(\overline{F(\zeta)},-1/\bar w)
=\rho(\tilde f(\zeta,w))$, using $F(\bar\zeta)=\overline{F(\zeta)}$, i.e.\ reality of $F$ on $\RR$.
Covering: the point $(\zeta,w)\in Q$ is $q=x+y\,\eta(w)$ where $\zeta=x+iy$; writing
$F(\zeta)=\alpha(x,y)+i\beta(x,y)$, the image point $(F(\zeta),w)$ is
$\alpha+\beta\,\eta(w)$, which equals $f(q)$ by \eqref{eq:intrinsic} applied in the slice
$\CC_{\eta(w)}$.
\end{proof}

\begin{remark}[Relation to the Gentili--Salamon--Stoppato lift]\label{rem:GSS}
Gentili--Salamon--Stoppato \cite[Theorem~5.3 and eq.~(5.6)]{GSS2014} construct the twistor lift of
an arbitrary slice regular function as a holomorphic map into $\CP^3$ (equivalently, a holomorphic
curve in the Klein quadric satisfying a reality condition, \cite[Theorem~5.7]{GSS2014}); the ruled
geometry of such lifts is developed by Altavilla and Altavilla--Sarfatti
\cite{Altavilla2018,AltavillaSarfatti2019}, and the $PGL(2,\HH)$-orbit structure of lifts has
recently been studied in \cite{LiuMorenoShi2026}. For an \emph{intrinsic} $f$ their lift formula
degenerates precisely to the product form of Theorem~\ref{thm:lift} (in their affine coordinates,
$(u,v)\mapsto(u,g(v))$ with $g$ the stem function; substitute $h=0$, $\hat g=g$ in
\cite[eq.~(5.6)]{GSS2014}). What we add here is the systematic use of this product form: the
identification of its natural domain and target as the incidence quadric $(Q,\rho)$ with its two
rulings, the equivariance, and the divisor-level dichotomy of Theorem~\ref{thm:dichotomy} below. We
have not found these statements in the literature, but the ingredients are all close to
\cite{GSS2014,Altavilla2018,AltavillaSarfatti2019}, and we make no strong novelty claim.
\end{remark}

The intrinsic class thereby acquires a geometric characterization: \emph{the rotationally symmetric
function theory on $\HH$ is the theory of holomorphic, $\rho$-equivariant self-maps of the incidence
quadric of product type $F\times\id$} --- a commutative function theory in one complex variable,
fibered trivially over the sphere of complex structures. Closure under addition, multiplication and
composition is visible by inspection in this model.

\subsection{Maximality of the quadric}\label{subsec:maximality}

In what sense is $Q$ the natural home of the lift, rather than an ambient $\CP^3$-neighborhood?
Some care is needed to make this question well posed: $Q$ is a smooth hypersurface in $\CP^3$, so
\emph{any} holomorphic map defined on $Q$ extends holomorphically to a neighborhood --- take local
coordinates in which $Q$ is a coordinate hyperplane and extend independently of the transverse
variable. Unconstrained local extension is therefore trivially possible and trivially
uninteresting; the correct question fixes the twistor structure. Two precise statements result: a
fiberwise computation (the extension cannot respect the incidence structure off $Q$), and a global
rigidity theorem (holomorphic immersions of the closed quadric into $\CP^3$ are projectively
linear, so no non-M\"obius lift extends to an immersive map of $Q$ into the ambient space other
than through $Q$ itself).

\begin{proposition}[The na\"ive extension is holomorphic only along $Q$]\label{prop:naive}
Fix an intrinsic slice regular $f$ on an axially symmetric slice domain, not affine
($f\not\equiv aq+b$, $a,b\in\RR$), and consider the fiber-preserving assignment
\[
\Psi_f\colon\ \HH\times\SS\to\HH\times\SS,\qquad (q,\eta)\mapsto(f(q),\eta),
\]
on the full correspondence space, where the complex structure at $(q,\eta)$ is left multiplication
by $\eta$ on the $\HH$-factor and the standard one (via the coordinate $w$, $\eta=\eta(w)$) on the
$\SS$-factor. Then the Cauchy--Riemann system for $\Psi_f$ holds at every point of $Q$
(where it reduces to slice holomorphy), and there is \emph{no open subset of}
$(\HH\times\SS)\setminus Q$ on which it holds: the locus where $\Psi_f$ satisfies the
Cauchy--Riemann equations is $Q\cup E$ with $E$ a proper real-analytic subset of the complement,
empty for $f=q^2$.
\end{proposition}

\begin{proof}
Holomorphy at $(q,\eta)$ in the $\HH$-directions requires the real differential $d_qf$ to
intertwine left multiplication by $\eta$: $d_qf(\eta v)=\eta\,d_qf(v)$ for all $v\in\HH$.
Decompose $T_q\HH=\CC_J\oplus\CC_J^\perp$ where $q=x+yJ$, $y>0$. For a monomial $f=q^m$ one
computes, for $v$ in the slice plane, $d_qf(v)=m q^{m-1}v$ if $v\in\CC_J$ (slice holomorphy), while
for $v\perp\CC_J$, using $vq=\bar qv$,
\[
d_qf(v)=\Bigl(\sum_{i=0}^{m-1}q^i\bar q^{\,m-1-i}\Bigr)v=\frac{q^m-\bar q^m}{q-\bar q}\,v .
\]
(The middle expression is a two-sided sum collapsing to the divided difference because
$q\bar q=|q|^2$ is real; see also Lemma~\ref{lem:jacobian} below, where these two blocks reappear
as the slice and spherical derivatives.) By $\RR$-linearity and locally uniform convergence the
same block decomposition holds for any slice regular $f$:
\begin{equation}\label{eq:blocks}
d_qf(v)=f'(q)\,v \quad (v\in\CC_J),\qquad
d_qf(v)=\partial_s f(q)\,v \quad (v\in\CC_J^\perp),
\end{equation}
where $f'$ is the slice derivative and
$\partial_sf(q)=(f(q)-f(\bar q))(q-\bar q)^{-1}$ is the \emph{spherical derivative}
\cite{GhiloniPerotti2011,GSSbook}. Both act by left multiplication by elements of $\CC_J$. Now take
$\eta\notin\{\pm J\}$ and decompose $\eta=\eta_\parallel+\eta_\perp$ with
$\eta_\parallel\in\CC_J$, $\eta_\perp\in\CC_J^\perp$, $\eta_\perp\neq0$. Testing the intertwining
condition on $v\in\CC_J$: $\eta v=\eta_\parallel v+\eta_\perp v$ splits into slice and orthogonal
parts, so
\[
d_qf(\eta v)=f'(q)\eta_\parallel v+\partial_sf(q)\eta_\perp v,
\qquad
\eta\,d_qf(v)=\eta_\parallel f'(q)v+\eta_\perp f'(q)v .
\]
The slice parts agree ($\CC_J$ is commutative); the orthogonal parts require
$\partial_sf(q)\,\eta_\perp v=\eta_\perp f'(q)v$, i.e.\ (since $\eta_\perp w=\bar w\eta_\perp$ for
$w\in\CC_J$) the pointwise identity $\partial_sf(q)=\overline{f'(q)}$. For intrinsic $f$ with
stem $F=\alpha+i\beta$, the spherical derivative $\partial_sf(x+yJ)=\beta(x,y)/y$ is
\emph{real}, so on any open set where the identity holds, $F'$ is real-valued; a holomorphic
function with real values on an open set is constant (open mapping theorem), and a real constant
$F'\equiv a$ gives $F$ affine with real coefficients --- and then $\beta/y=a=F'$, so the identity
does hold, consistently with the M\"obius case. Hence for non-affine intrinsic $f$ the
Cauchy--Riemann system fails on a dense open subset of every neighborhood in
$(\HH\times\SS)\setminus Q$; the exceptional locus $E$ where the pointwise identity happens to
hold off $Q$ is the real-analytic set $\{\partial_sf=\overline{f'}\}$, proper by the above (and
empty off $Q$ for $f=q^2$: there the identity reads $2x=2\bar q$, forcing $q\in\RR$, which lies
on $Q$). On $Q$ itself ($\eta=\pm J$; the copy $\eta=-J$ is
identified with $Q$ under the antipode $w\mapsto-1/\bar w$), the intertwining condition is exactly
slice holomorphy, which holds. The fiber directions impose nothing: $\Psi_f$ is the identity in
$\eta$ at fixed $q$, holomorphic for the standard structure.
\end{proof}

\begin{proposition}[Rigidity of quadric immersions]\label{prop:immersion}
Every holomorphic immersion $\Phi\colon\CP^1\times\CP^1\to\CP^3$ is of the form
$\Phi=p\circ\iota\circ(m\times\tilde m)$, possibly precomposed with the factor swap, where $\iota$
is the Segre embedding, $p\in PGL_4(\CC)$, and $m,\tilde m\in PGL_2(\CC)$. In particular the image
is a smooth quadric, and $\Aut(Q)\cong(PGL_2(\CC)\times PGL_2(\CC))\rtimes\ZZ/2$ consists entirely
of restrictions of projective transformations of\, $\CP^3$.
\end{proposition}

\begin{proof}
Write $S=\CP^1\times\CP^1$ and $L=\Phi^*\mathcal O(1)=\mathcal O(a,b)$. The image spans $\CP^3$:
if it lay in a hyperplane $\CP^2$, then $\Phi$ would be a holomorphic immersion of a compact
surface into a surface, hence a covering map onto the simply connected $\CP^2$, forcing
$S\cong\CP^2$, false. If $a=0$ (or $b=0$), $\Phi$
is constant on the corresponding ruling family in one direction and cannot be immersive, so
$a,b\ge1$. The line bundle $\mathcal O(a,b)$ with $a,b\ge1$ is very ample with complete linear
system of dimension $N=(a+1)(b+1)-1$, embedding $S$ into $\CP^N$; $\Phi$ is the composition of
this embedding with a linear projection $\CP^N\dashrightarrow\CP^3$ from a center
$\Lambda\cong\CP^{N-4}$ disjoint from $S$. Such a projection is an immersion if and only if
$\Lambda$ avoids the embedded tangent variety $\mathrm{Tan}(S)\subset\CP^N$. Since $S$ is smooth,
non-planar and not a cone or a developable, $\dim\mathrm{Tan}(S)=4$ (surfaces with degenerate
tangent varieties are classified; see \cite{GriffithsHarris1979} or \cite[Ch.~I]{Zak1993}).
As $\dim\Lambda+\dim\mathrm{Tan}(S)=N-4+4=N$, in $\CP^N$ the center always meets
$\mathrm{Tan}(S)$ unless the projection is trivial, i.e.\ unless $N=3$. Hence $(a,b)=(1,1)$,
$\Phi$ is the Segre embedding composed with $p\in PGL_4$, and the identifications of the factors
are M\"obius. For the last statement: any automorphism of $Q$ preserves
$-K_Q=\mathcal O(2,2)$ and hence its unique square root $\mathcal O(1,1)$ (the Picard group is
torsion free), and $Q\subset\CP^3$ is embedded by the complete system $|\mathcal O(1,1)|$, so the
automorphism extends to $\CP^3$; explicitly, $\CP^3=\mathbb P(\CC^2\otimes\CC^2)$, both $PGL_2$
factors act through the tensor product and the ruling swap is induced by $u\otimes v\mapsto
v\otimes u$. The classical background on quadrics is in \cite[Ch.~6]{GriffithsHarris1978}.
\end{proof}

\begin{remark}[What maximality does and does not say]\label{rem:maximality}
Proposition~\ref{prop:naive} says the incidence relation $q\in\CC_\eta$ is forced: the lift cannot
be prolonged off $Q$ compatibly with the fibration by complex structures unless $f$ is M\"obius
(over $\RR$, affine; adding the point at infinity, fractional linear). Proposition~\ref{prop:immersion}
complements it globally: the closed quadric admits no immersive holomorphic repositioning in
$\CP^3$ other than the linear ones, so there is no ``bigger'' holomorphic model of the intrinsic
class inside twistor space --- and by Riemann--Hurwitz, $\tilde f=F\times\id$ composed with the
Segre embedding is an immersion if and only if $F$ has no critical points on $\CP^1$, i.e.\ if and
only if $F$ is M\"obius. This resolves the rigidity tension of \cite[Prop.~3.1]{Baker2017} (a
holomorphic self-map of twistor space restricting to degree-one maps on all fibers is a M\"obius
transformation): the rotational function class does not live on twistor space; it lives on the
incidence quadric. We do not claim, and do not need, a statement excluding arbitrary
(fibration-incompatible) local extensions, which exist trivially across any smooth hypersurface.
The na\"ive-extension computation for $f=q^2$ appears in the working notes \cite{ZenodoArchive};
the general case is Proposition~\ref{prop:naive}. A version of Proposition~\ref{prop:immersion}
may well be classical; we could not locate a reference stating it in this form, and the proof
above is assembled from classical facts. (We take this occasion to note that the smooth-topology
paper of Li--Peterson on immersions of $n$-manifolds in $(2n{-}2)$-manifolds, sometimes cited in
adjacent contexts, concerns a different category and does not apply here.)
\end{remark}

\subsection{The equivariant dichotomy}

Consider now the zero divisor of the lifted function. For intrinsic $f$ with stem $F$, the zero
divisor of $\tilde f$ in $Q$ is
\[
Z(\tilde f)=F^{-1}(0)\times\CP^1_w=\sum_{F(\zeta_0)=0}L_{\zeta_0},
\]
a union of ruling-A lines with multiplicities, one per zero of $F$.

\begin{example}\label{ex:modelcases}
$f(q)=q^2-1$: $F(\zeta)=\zeta^2-1$, zeros $\pm1$ real; $Z(\tilde f)=L_1+L_{-1}$ is a pair of
$\sigma$-fixed lines (twistor fibers); downstairs, two real point zeros.
$f(q)=q^2+1$: $F(\zeta)=\zeta^2+1$, zeros $\pm i$; $Z(\tilde f)=L_i+L_{-i}$ is a
$\sigma$-\emph{swapped pair} of lines, each individually not invariant, each transverse to the
fibers; downstairs, the single zero sphere $\SS$.
\end{example}

In both cases the homology class of the divisor in $H_2(Q;\ZZ)$ is $2[\text{ruling A}]$ ---
identical. The discriminating invariant is equivariant, not homological:

\begin{mainthm}[Equivariant dichotomy]\label{thm:dichotomy}
Let $f\not\equiv0$ be intrinsic slice regular on an axially symmetric slice domain, with lift
$\tilde f=F\times\id$ and zero divisor $Z(\tilde f)=\sum_k m_k L_{\zeta_k}$. Then:
\begin{enumerate}
\item Each component $L_{\zeta_k}$ is a ruling-A line; $\sigma(L_{\zeta_k})=L_{\bar\zeta_k}$.
A real zero of $f$ contributes a $\sigma$-\emph{fixed} component (a twistor fiber); a zero sphere
contributes a $\sigma$-\emph{swapped pair} of components. The multiset $\{\zeta_k\}$ with
multiplicities is $\sigma$-invariant.
\item In $H_2(Q;\ZZ)$, a real zero of multiplicity $m$ contributes $m[\mathrm{ruling~A}]$ and a
zero sphere of multiplicity $m$ contributes $2m[\mathrm{ruling~A}]$ (one per swapped component);
in particular a pair of simple real zeros and a single simple sphere have the same class
$2[\mathrm{ruling~A}]$: homology cannot separate the two cases.
\item The quotient $Y=Q/\sigma$ is a smooth $4$-manifold ($\sigma$ acts freely), fibered over the
closed half-plane $\{\im\zeta\ge0\}\cup\{\infty\}$; a real zero descends to a component $\cong\RP^2$
(nonorientable) and a zero sphere to a component $\cong S^2$ (orientable). Consequently, for
$f=\Xi$:
\[
\text{RH}\iff\text{every component of the zero locus of }\tilde\Xi\text{ in }Y=Q/\sigma
\text{ is nonorientable.}
\]
\end{enumerate}
\end{mainthm}

\begin{proof}
(1) The divisor description precedes the statement; $\sigma(L_\zeta)=L_{\bar\zeta}$ was computed
after \eqref{eq:line}; the multiset is $\sigma$-invariant because $F$ is real on $\RR$, so its
zeros are closed under conjugation with multiplicity.

(2) In $H_2(Q;\ZZ)=\ZZ[\text{A}]\oplus\ZZ[\text{B}]$ each line $L_\zeta$ has class $[\text{A}]$. A
real zero of the intrinsic $f$ at $x_0$ of multiplicity $m$ downstairs corresponds to
$F(\zeta)=(\zeta-x_0)^m(\cdots)$, contributing $m[\text{A}]$; a sphere $(a,y_0)$ of multiplicity
$m$ contributes the pair $mL_{a+iy_0}+mL_{a-iy_0}$, class $2m[\text{A}]$. Two simple real zeros and
one simple sphere both give $2[\text{A}]$.

(3) $\sigma$ has no fixed points on $\CP^3$ (hence none on $Q$), so $Y$ is a smooth compact
$4$-manifold and $Q\to Y$ is a double cover. In coordinates \eqref{eq:rho}, $\rho$ covers
$\zeta\mapsto\bar\zeta$ on the first factor, so $Y$ fibers over
$\{\im\zeta\ge0\}\cup\{\infty\}$; over interior points the fiber is the $2$-sphere $\CP^1_w$ (the
two lines $L_\zeta,L_{\bar\zeta}$ are identified into one), while over boundary points
($\zeta\in\RR\cup\{\infty\}$) the fiber is $L_\zeta/(\text{antipodal})\cong\RP^2$. A
$\sigma$-fixed component $L_{x_0}$ descends to the boundary fiber $\RP^2$; a $\sigma$-swapped pair
$L_{a\pm iy_0}$ descends to a single interior fiber $S^2$. The displayed equivalence is then
Theorem~\ref{thm:reformulation} read in $Y$.
\end{proof}

\begin{remark}[Value-curve form]\label{rem:valuecurve}
An equivalent formulation, useful for enumerative experiments: the Pl\"ucker image of the moving
line $\zeta\mapsto L_{F(\zeta)}$ is the curve $[F^2:0:F:-F:0:1]$ in the Klein quadric, which
passes through the fiber point $[0:\cdots:0:1]$ once per zero of $F$. At a double passage the
curve has a node whose two branches are individually $\sigma$-real for a pair of real zeros
(crunode type), and exchanged by $\sigma$ for a conjugate pair (acnode/solitary type). Real zeros
versus zero spheres is the hyperbolic-versus-elliptic dichotomy of real nodes of the value curve
at real points of the Klein quadric. Under the de Bruijn--Newman heat flow (\S\ref{sec:hardness}),
an acnode converts to a crunode through a tacnode --- in $Y$, an $S^2$ component collides with the
$\RP^2$ boundary stratum and splits into two $\RP^2$'s.
\end{remark}

\begin{remark}[Positioning]\label{rem:positioning}
The lift-and-divisor picture should be read against three bodies of work. The twistor lifts of
slice regular functions and the induced orthogonal complex structures (OCS) are due to
Gentili--Salamon--Stoppato \cite{GSS2014}; the ruled-surface geometry of such lifts was developed
by Altavilla and Altavilla--Sarfatti \cite{Altavilla2018,AltavillaSarfatti2019}; and the rigidity
and removability theory for OCS on domains of $\RR^4$, which constrains what the lift geometry can
do at singularities, is Salamon--Viaclovsky \cite{SalamonViaclovsky2009} (in particular their
Liouville-type Theorem~1.3 and the quadric constructions of their Theorem~1.13). The equivariant
component-type invariant of Theorem~\ref{thm:dichotomy}, elementary as it is, does not appear in
these works to our knowledge; its interest is that it is precisely \emph{not} a deformation
invariant (components change type exactly when zeros collide on the axis), which by
Theorem~\ref{thm:degree-nogo} and the hardness checks of Section~\ref{sec:hardness} is a necessary
feature of any true discriminant.
\end{remark}

%=====================================================================
\section{Two no-go theorems}\label{sec:nogo}

The reformulation invites a search for a functional that ``sees'' spheres. This section proves
that the two most natural classes of candidates --- linear functionals of rotational extensions,
and deformation-invariant degree counts --- cannot work. Both results are dictated by the hardness
checks: the Davenport--Heilbronn function (\S\ref{sec:hardness}) shows that soft mechanisms must
fail, and these theorems say \emph{how} they fail.

\subsection{The linear no-go}

The slice extension $R\colon F\mapsto f$ of \eqref{eq:intrinsic} and the Fueter-side hierarchy
$Q_\ell=(B_+)^\ell\circ Q_0$ of \cite{Baker2015} (recalled in \S\ref{sec:sectors}) are both
$\RR$-linear in the complex input and injective on the relevant classes. This has a blunt
consequence.

\begin{mainthm}[Linear no-go]\label{thm:linear-nogo}
Let $\mathcal O_{\mathrm{int}}(\CC)$ denote the intrinsic entire functions (real Taylor
coefficients), with the Fr\'echet topology of locally uniform convergence, and let $E$ be any
$\RR$-linear continuous ``rotational extension'' map from $\mathcal O_{\mathrm{int}}(\CC)$ to a
topological $\RR$-vector space of $\HH$-valued functions on $\HH$ (e.g.\ the slice extension
$E=R$, or the Fueter-side maps $E=Q_\ell$ of \S\ref{sec:sectors}, any $\ell$). Then:
\begin{enumerate}
\item every continuous $\RR$-linear functional $\Lambda$ of $E(F)$ is a continuous
$\RR$-linear functional $\Lambda_0=\Lambda\circ E$ of $F$ itself: the extension carries no new
linearly accessible information;
\item no continuous $\RR$-linear functional detects reality of zeros: there exist no continuous
$\RR$-linear $\Lambda_0$ on $\mathcal O_{\mathrm{int}}(\CC)$ and threshold $c\in\RR$ such that,
on any subset of $\mathcal O_{\mathrm{int}}(\CC)$ containing the two-parameter family
$F_{t,s}(\zeta)=(\zeta^2+t)(\zeta^2+s)$, the reality locus
$\{F:\text{all zeros of }F\text{ real}\}$ coincides with $\{\Lambda_0(F)\ge c\}$ or with
$\{\Lambda_0(F)=c\}$.
\end{enumerate}
In particular, no linear pairing of any rotational extension of\, $\Xi$-like functions ---
against any fixed test function, over any contour, in any sector --- can separate on-axis from
off-axis zeros.
\end{mainthm}

\begin{proof}
(1) is the composition of continuous $\RR$-linear maps.

(2) The zeros of $F_{t,s}$ are all real if and only if $t\le0$ and $s\le0$ (a quarter-plane
$P$ in the $(t,s)$-parameters). Expanding,
$F_{t,s}=\zeta^4+(t+s)\zeta^2+ts$, so by linearity
\[
\Lambda_0(F_{t,s})=A+(t+s)B+ts\,C,\qquad
A=\Lambda_0(\zeta^4),\ B=\Lambda_0(\zeta^2),\ C=\Lambda_0(1).
\]
Suppose the reality locus equals $\{\Lambda_0\ge c\}$. The boundary of $P$ contains the ray
$\{t=0,\,s\le0\}$, along which $\Lambda_0(F_{0,s})=A+sB$ must equal the threshold behavior of a
boundary: since points with $t=0^+$ are outside $P$ and $t=0^-$ inside, continuity forces
$A+sB=c$ for all $s\le0$, hence $B=0$ and $A=c$. The condition $\Lambda_0\ge c$ then reads
$ts\,C\ge0$, whose solution set is $\{ts\ge0\}$ (if $C>0$), $\{ts\le0\}$ (if $C<0$), or
everything (if $C=0$) --- none of which equals the quarter-plane $P$. The case
$\{\Lambda_0=c\}$ is identical: the level set $\{A+(t+s)B+tsC=c\}$ is empty, a line, all of
$\RR^2$, a hyperbola, or a pair of crossed full lines --- never the boundary of a quarter-plane
(two half-lines meeting at a corner).
\end{proof}

\begin{remark}\label{rem:nogo-history}
Theorem~\ref{thm:linear-nogo} formalizes and generalizes the negative result of
\cite[\S4]{Baker2015}, where it was proved by an explicit Laurent-expansion computation that the
$4$-dimensional pairings of the Fueter extension against rotationally symmetric regular test
functions reproduce the $2$-dimensional contour integrals of the complex input exactly ---
including for the kernels singular on spheres. That computation is the concrete instance
($E=Q_0$, $\Lambda=$ Cauchy--Fueter pairing) of the abstract statement above; the theorem explains
why a decade of linear attempts had to fail and redirects the search to functionals that are
nonlinear (quadratic, \S\ref{sec:positivity}--\ref{sec:bridge}), geometric-equivariant
(\S\ref{sec:twistor}), or spectral.
\end{remark}

\subsection{The degree no-go}

The second candidate is topological: the quaternionic winding number, i.e.\ the Brouwer degree of
$f/|f|$ on $3$-cycles, which is the natural $4$-dimensional argument principle. We first record
the Jacobian structure of intrinsic maps.

\begin{lemma}[Jacobian factorization]\label{lem:jacobian}
Let $f$ be slice regular near $q=x+yJ$, $y>0$. With respect to the orthogonal splitting
$T_q\HH=\CC_J\oplus\CC_J^\perp$, the real differential of $f$ acts as left multiplication by
$f'(q)$ on $\CC_J$ and by the spherical derivative $\partial_sf(q)=(f(q)-f(\bar q))(q-\bar q)^{-1}$
on $\CC_J^\perp$, as in \eqref{eq:blocks}. Consequently
\[
\det(d_qf)=|f'(q)|^2\,|\partial_sf(q)|^2\ \ge 0 :
\]
slice regular maps are orientation-preserving away from the degenerate set
$\{f'=0\}\cup\{\partial_sf=0\}$.
\end{lemma}

\begin{proof}
The block decomposition \eqref{eq:blocks} was established in the proof of
Proposition~\ref{prop:naive} (for monomials by the divided-difference computation, extended by
linearity and locally uniform convergence; the identification of the orthogonal block with the
spherical derivative is the standard one \cite{GhiloniPerotti2011,GSSbook}). Left multiplication
by $a\in\HH$ on a $2$-plane it preserves acts with determinant $|a|^2$; both blocks are left
multiplications by elements of $\CC_J$, each preserving its plane, giving the product formula.
\end{proof}

\begin{mainthm}[Degree no-go]\label{thm:degree-nogo}
Let $f$ be intrinsic slice regular. The local Brouwer degree of $f/|f|\colon\partial N\to S^3$ on
the boundary of a small tubular neighborhood $N$ of a zero component is: $m$ at a real zero of
multiplicity $m$, and $2m$ at a zero sphere of multiplicity $m$. That is, the quaternionic winding
number weights a zero sphere exactly as it weights the corresponding conjugate pair of complex
zeros: deformation-invariant integer counts cannot detect spheres.
\end{mainthm}

\begin{proof}
\emph{Real zero.} For $f=q^m$ on $|q|=1$, a generic target value has exactly the $m$ slice roots
as preimages, and by Lemma~\ref{lem:jacobian} each preimage is positively oriented (on the slice
plane $\det=m^2|q|^{2(m-1)}$; on the orthogonal plane the divided difference
$(q^m-\bar q^m)(q-\bar q)^{-1}=\sin m\psi/\sin\psi\cdot$(unit) for $q=e^{J\psi}$, again of
nonnegative determinant, and positive at generic points); hence local degree $m$. A general
intrinsic $f$ with a real zero of multiplicity $m$ factors locally as $(q-x_0)^m\tilde f(q)$ with
$\tilde f(x_0)\neq0$ real, and the degree is stable under this real perturbation.

\emph{Zero sphere.} For the model $f=q^2+1$ near the sphere $\SS$: writing $q=x+yJ$, the local
model is $f(x+yJ)=(x^2-y^2+1)+2xyJ\approx-2(y-1)+2xJ$ near $(x,y)=(0,1)$. On the boundary
$S^1_\phi\times S^2_J$ of a tubular neighborhood of the zero sphere the map is
$(\phi,J)\mapsto\cos\phi'+J\sin\phi'$ (with $\phi'$ the angular coordinate of
$(-2(y-1),2x)$), which hits a generic target value twice, at $(\phi_0,J_0)$ and $(-\phi_0,-J_0)$.
Both preimages count positively: by Lemma~\ref{lem:jacobian} the map is orientation-preserving
where nondegenerate, and global consistency forces the count. Indeed the family $q^2+t$,
$t\in[0,1]$, is nonvanishing on the fixed large sphere $|q|=R>1$, so the degree there is constant
in $t$; at $t=0$ it is the degree of $q^2$, which equals $2$ by the real-zero computation above
(a real zero of multiplicity $2$); and at $t=1$ the zero set is the unit sphere $\SS$. In
$H_3(\HH\setminus\SS)\cong\ZZ$ the large sphere $S^3_R$ is homologous to the boundary
$S^2\times S^1$ of the tubular neighborhood (the region between them is a $4$-chain avoiding
$\SS$), so the local degree at the zero sphere is $2$. Multiplicity
$m$ gives $2m$ by the same argument applied to $(q^2+1)^m\cdot(\text{unit})$.
\end{proof}

\begin{remark}\label{rem:c2}
The count $r+2s$ (real zeros plus twice spheres) is exactly the count the complex argument
principle produces for the even function $\Xi$ --- the ``$c=2$ conspiracy'' of
\cite[\S7]{Baker2015} persists for the nonlinear degree. This is consistent with, and required
by, the hardness checks: the degree is a deformation invariant, spheres are destroyed by
arbitrarily small deformations (\S\ref{sec:hardness}), so the degree \emph{must not} see them.
Any successful discriminant must be finer than homology and finer than degree; the equivariant
component type of Theorem~\ref{thm:dichotomy} is the minimal such refinement, and it is not
deformation-invariant --- components change type precisely when zeros collide on the real axis.
\end{remark}

%=====================================================================
\section{Hardness checks}\label{sec:hardness}

Any proposed sphere-excluding mechanism must pass two tests. We state them as requirements,
with the data that makes them concrete.

\subsection{The Davenport--Heilbronn test}\label{subsec:DH}

Let $\chi$ be the odd Dirichlet character mod $5$ with $\chi(2)=i$, and set
\[
\kappa=\frac{\sqrt{10-2\sqrt5}-2}{\sqrt5-1},\qquad
f_{DH}(s)=\tfrac12\bigl[(1-i\kappa)L(s,\chi)+(1+i\kappa)L(s,\bar\chi)\bigr],
\]
\[
\Lambda_{DH}(s)=\Bigl(\frac{5}{\pi}\Bigr)^{s/2}\Gamma\Bigl(\frac{s+1}{2}\Bigr)f_{DH}(s),
\qquad \Xi_{DH}(z)=\Lambda_{DH}\bigl(\tfrac12+iz\bigr).
\]
Then $\Xi_{DH}$ is entire of order $1$, even, and real on $\RR$ (evenness and reality verified
numerically to relative $10^{-16}$ in Appendix~\ref{app:numerics}, and provable directly from
the functional equations of $L(s,\chi)$, $L(s,\bar\chi)$ and the choice of $\kappa$), so it
possesses an intrinsic slice regular extension
with exactly the symmetry of $\Xi$ --- but $f_{DH}$ has zeros off the critical line
\cite{DavenportHeilbronn1936a,DavenportHeilbronn1936b,Titchmarsh1986}, and no Euler product. Its
quaternionic extension therefore has genuine zero spheres. The first three, recomputed here by
Newton iteration at 25-digit precision (Appendix~\ref{app:numerics}) and matching Spira's
computations \cite{Spira1994} (see also \cite{BalanzarioSanchez2007}):
\begin{equation}\label{eq:DHspheres}
\begin{array}{lll}
s=0.808517182\ldots+85.699348485\ldots i, & a=85.699348, & y_0=0.308517,\\
s=0.650830081\ldots+114.163342686\ldots i, & a=114.163343, & y_0=0.150830,\\
s=0.574356050\ldots+166.479305\ldots i, & a=166.479306, & y_0=0.074356 .
\end{array}
\end{equation}
(Here $(a,y_0)$ are the center and radius of the corresponding sphere in rotated coordinates:
$s=\sigma+i\tau$ off the line gives $\gamma=\tau- i(\sigma-\tfrac12)$ and the sphere
$a=\tau$, $y_0=\sigma-\tfrac12$.)

\textbf{Requirement 1.} \emph{Symmetry plus growth cannot exclude spheres; a candidate mechanism
must consume arithmetic input (Euler product / explicit formula / positivity of $\Phi$ used in an
essential way), and must fail --- demonstrably --- for $\Xi_{DH}$.} Every structure proposed in
this paper is run against DH: the Laguerre detector dips negative exactly inside the predicted
windows at the spheres \eqref{eq:DHspheres} (\S\ref{sec:positivity}), the sector energies pick up
exactly the predicted $y_0^{2\ell}$ signal (\S\ref{sec:positivity}), and the growth criterion
recovers the first radius to four decimals (\S\ref{sec:bridge}).

\subsection{The small-sphere (de Bruijn--Newman) test}\label{subsec:dBN}

For $t\in\RR$ let $H_t(z)=\int_0^\infty e^{tu^2}\,\Phi(u)\cos(zu)\dd u$ (up to normalization,
the heat-flow deformation of $\Xi$), and let $\Lambda_{dBN}$ be the de Bruijn--Newman constant:
$H_t$ has only real zeros iff $t\ge\Lambda_{dBN}$, and RH is equivalent to $\Lambda_{dBN}\le0$.
Rodgers--Tao \cite{RodgersTao2020} proved $\Lambda_{dBN}\ge0$: RH, if true, is \emph{barely}
true. In the present language: for $t<0$ the deformations $H_t$ have zero spheres of arbitrarily
small radius accumulating along the axis, and under the flow a sphere shrinks, collides with the
axis, and splits into two real zeros (the tacnode transition of Remark~\ref{rem:valuecurve});
$\Xi$ sits exactly at the closure of the sphere-free regime.

\textbf{Requirement 2.} \emph{Any true criterion degrades only polynomially or worse as the sphere
radius $y_0\to0$; a soft argument stable under small perturbations of the function proves too much
and is therefore wrong.} The structures of \S\S\ref{sec:positivity}--\ref{sec:bridge} comply
quantitatively: the sector-$\ell$ signal of a sphere scales as $y_0^{2\ell}$
(Lemma~\ref{lem:grading}), and the growth criterion needs frequencies $\lambda\gtrsim1/y_0$
(Theorem~\ref{thm:rate} and Remark~\ref{rem:honest}) --- arbitrarily weak, never zero.

%=====================================================================
\section{The positivity hierarchy}\label{sec:positivity}

The no-go theorems force nonlinearity; the simplest nonlinear invariants are quadratic. This
section develops the quadratic family on the ``unitary'' side: a hierarchy of positivity
statements graded by spherical Bessel tags, whose bottom rung is the classical Laguerre
inequality. Section~\ref{sec:sectors} grounds the tags in the quaternionic harmonic analysis of
\cite{Baker2015,Baker2016}, and Section~\ref{sec:bridge} develops the ``analytic'' twin.

\subsection{The bottom rung: Laguerre positivity}\label{subsec:laguerre}

For $\Xi$ real entire of order $1$ with Hadamard product over zeros $\rho$ (in the rotated
variable, so RH says $\rho\in\RR$),
\begin{equation}\label{eq:logsecond}
-(\log\Xi)''(x)=\sum_\rho\frac{1}{(x-\rho)^2},\qquad x\in\RR,\ \Xi(x)\neq0 .
\end{equation}
A real zero contributes $(x-\rho)^{-2}>0$; a conjugate pair $\rho=a\pm iy_0$ (a zero sphere of
center $a$, radius $y_0$) contributes
\[
2\,\re\frac{1}{(x-a-iy_0)^2}
=\frac{2\bigl((x-a)^2-y_0^2\bigr)}{\bigl((x-a)^2+y_0^2\bigr)^2},
\]
which is \emph{negative exactly on the interval} $|x-a|<y_0$ --- the equatorial shadow of the
sphere. Positivity of \eqref{eq:logsecond} against $\Xi^2$ is the Laguerre inequality
\[
L_1[\Xi](x)=\Xi'(x)^2-\Xi(x)\Xi''(x)\ge0\qquad(x\in\RR),
\]
a necessary condition for RH. The full hierarchy characterizes RH:

\begin{theorem}[Csordas--Varga; Csordas--Vishnyakova]\label{thm:laguerre}
For a real entire function $\varphi\not\equiv0$ define
\[
L_n(x;\varphi)=\sum_{j=0}^{2n}\frac{(-1)^{j+n}}{(2n)!}\binom{2n}{j}
\varphi^{(j)}(x)\,\varphi^{(2n-j)}(x),
\qquad\text{so that }\ |\varphi(x+iy)|^2=\sum_{n\ge0}L_n(x;\varphi)\,y^{2n}.
\]
If $\varphi$ belongs to the Laguerre--P\'olya class then $L_n(x;\varphi)\ge0$ for all $n\ge0$ and
all $x\in\RR$ \textup{(\cite[Theorem~2.9]{CsordasVarga1990})}. Conversely, if
$L_n(x;\varphi)\ge0$ for all $n\ge0$ and all $x\in\RR$, then $\varphi$ belongs to the
Laguerre--P\'olya class \textup{(\cite[Theorem~2.3]{CsordasVishnyakova2013}; no a priori
order or genus restriction)}. Since $\Xi$ is real entire of order $1$, membership
$\Xi\in\LP$ is equivalent to all zeros of $\Xi$ being real, so
\[
\text{RH}\iff L_n(x;\Xi)\ge0\ \text{ for all }n\in\NN_0\text{ and all }x\in\RR .
\]
\end{theorem}

\begin{remark}
For $\Xi$ the ``inequalities $\Rightarrow$ reality'' direction is also elementary, without
citation: if $\Xi(x_0+iy_0)=0$ with $y_0\neq0$, then
$0=|\Xi(x_0+iy_0)|^2=\sum_nL_n(x_0;\Xi)\,y_0^{2n}$; if all $L_n(x_0)\ge0$ then every term
vanishes, so $|\Xi(x_0+iy)|=0$ for all $y$, forcing $\Xi\equiv0$.
\end{remark}

In the twistor language of Section~\ref{sec:twistor}, $L_1\ge0$ is the statement that the value
curve meets the real Klein quadric ``from the crunode side'' at each passage
(Remark~\ref{rem:valuecurve}); the higher $L_n$ are its jets, cf.\ the expansion of
$|\varphi(x+iy)|^2$ in the sphere-radius variable $y$. The hyperbolicity of Jensen polynomials
--- equivalent to RH in the aggregate, and proved for all but finitely many polynomials of each
degree by Griffin--Ono--Rolen--Zagier \cite{GORZ2019} --- is the polynomial-truncation face of
the same Laguerre--P\'olya membership. The hierarchy complies with both hardness
checks: it fails for DH (next paragraph), and the sphere signal degrades polynomially in $y_0$.

\textbf{Numerics} (Appendix~\ref{app:numerics}, Experiments 1--2). On $[0,52]$, computed purely
from the cosine transform \eqref{eq:cosine}, the relative margin
$L_1[\Xi]/(\Xi'^2+|\Xi\Xi''|)$ has minimum $0.173$. For DH (after the sign-preserving rescaling
$W=e^{\pi x/4}\Xi_{DH}$, which leaves $\operatorname{sign}L_1$ invariant since
$L_1[e^{cx}g]=e^{2cx}L_1[g]$): $L_1[W]<0$ exactly on $[85.417,85.982]$ at the first sphere
(predicted window $|x-a|<y_0$: $[85.391,86.008]$) and on $[114.025,114.302]$ at the second
(predicted $[114.013,114.314]$) --- the observed windows sit just inside the predictions, the
shaving accounted for by the positive contributions of neighboring on-line zeros --- while on the
sphere-free control stretch $[78,82]$ the minimum is $\approx7.2\times10^2>0$ (rescaled units).
A zero sphere of radius $y_0$ forces $L_1<0$ on an interval of length $\approx2y_0$ at its
projection: \emph{the equivariant dichotomy of Theorem~\ref{thm:dichotomy} has an axis-visible,
quantitative certificate, computed from function data alone.}

\subsection{The graded family: sector energies}\label{subsec:sectors-energies}

The Laguerre inequality mixes all sphere radii. The grading that separates them is by spherical
Bessel tags. For a zero multiset $\{\rho\}$ written as $\rho=a_\rho+iy_\rho$ ($y_\rho\ge0$; each
conjugate pair counted once by its upper member), define the \emph{sector energies}
\begin{equation}\label{eq:sector}
\EE_\ell=\int_0^\infty\Bigl|\sum_\rho j_\ell(\lambda y_\rho)\,e^{i\lambda a_\rho}\Bigr|^2
w(\lambda)\dd\lambda,\qquad \ell=0,1,2,\ldots
\end{equation}
for a positive weight $w$ (the laboratory uses a Gaussian window), where $j_\ell$ is the
spherical Bessel function. These are quadratic in the zeros --- Theorem~\ref{thm:linear-nogo}
does not apply --- and graded:

\begin{lemma}[Grading]\label{lem:grading}
Let the zero multiset be finite (in practice: windowed) and closed under conjugation, and let
$w>0$ be continuous with all moments finite. Then:
\begin{enumerate}
\item Real zeros contribute nothing to any $\ell\ge1$ sector: $j_\ell(0)=0$ for $\ell\ge1$, so
if all $y_\rho=0$ then $\EE_\ell=0$ for every $\ell\ge1$, exactly.
\item A single sphere of radius $y_0$ contributes to $\EE_\ell$ at order $y_0^{2\ell}$ as
$y_0\to0$: from $j_\ell(t)=t^\ell/(2\ell+1)!!\,\bigl(1+O(t^2)\bigr)$,
\[
\EE_\ell(\text{one sphere }(a,y_0))
=\frac{y_0^{2\ell}}{((2\ell+1)!!)^2}\int_0^\infty\lambda^{2\ell}w(\lambda)\dd\lambda
\;\bigl(1+O(y_0^2)\bigr).
\]
\item All zeros are real $\iff$ $\EE_\ell=0$ for all $\ell\ge1$ $\iff$ $\EE_1=0$: the $\ell$-th
sector detects a sphere of radius $y_0$ with sensitivity $y_0^{2\ell}$, and the $\ell\ge1$
sectors jointly (indeed singly, at $\ell=1$) certify reality.
\end{enumerate}
\end{lemma}

\begin{proof}
(1) and (2) are immediate from the Taylor expansion of $j_\ell$ at $0$
(e.g.\ \cite[10.52.1]{DLMF}). For (3): suppose $\EE_1=0$. Since $w>0$ and the integrand is
continuous in $\lambda$, the finite sum $S(\lambda)=\sum_\rho j_1(\lambda y_\rho)e^{i\lambda
a_\rho}$ vanishes identically on $\lambda\ge0$. As $\lambda\to0^+$,
$j_1(\lambda y_\rho)=\tfrac13\lambda y_\rho+O(\lambda^3)$ and $e^{i\lambda a_\rho}=1+O(\lambda)$,
so $S(\lambda)=\tfrac{\lambda}{3}\sum_\rho y_\rho+O(\lambda^2)$; hence $\sum_\rho y_\rho=0$,
and since every $y_\rho\ge0$, all vanish.
\end{proof}

\textbf{Numerics} (Appendix~\ref{app:numerics}, Experiment 3). With $20$ real Riemann ordinates
plus one implanted sphere at $a=28$ of radius $y_0$ swept over $[10^{-3},10^{-0.5}]$, and
Gaussian weight: the measured log--log slopes of $\EE_\ell$ versus $y_0$ are $2.000$, $4.000$,
$6.000$ for $\ell=1,2,3$; $\EE_1$ of the real ordinates alone is exactly $0$; and on real DH
data the three spheres \eqref{eq:DHspheres} give per-sphere $\EE_1$ values
$4.360\times10^{-3}$, $1.062\times10^{-3}$, $2.593\times10^{-4}$, with successive ratios
$4.11$, $4.10$ against the predicted $(y_0/y_0')^2=4.18$, $4.11$.

The polynomial (never exponential, never zero) decay in $y_0$ is the quantitative face of
Requirement 2: near-axis spheres emit an arbitrarily weak but nonvanishing signal, so a
successful positivity proof must control these functionals to all orders --- which is where the
arithmetic must enter, since DH shows that symmetry alone permits the signal to be present.

\subsection{Where the tags come from: the quaternionic sector decomposition}\label{sec:sectors}

The Bessel tags in \eqref{eq:sector} are not ad hoc: they are the radial profiles of the
$SU(2)$-isotypic decomposition of the rotationally extended function theory, worked out in
\cite{Baker2015,Baker2016}, which we now summarize in the form needed. (This subsection uses the
Fueter-regular extension hierarchy; it is the harmonic-analysis complement of the slice-regular
algebra used elsewhere, the two being linked by Fueter's theorem. See
\cite{BoryReyesShapiro2010} for the dictionary between the quaternionic and Clifford-analytic
settings, in which this machinery has parallel developments
\cite{BrackxDeSchepperSommen2005,DeBieXu2011}.)

Rotations about the real axis act on $\HH$-valued functions by
$M_\eta(F)(q)=\bar\eta F(\eta q\bar\eta)$, $\eta\in\HH$, $|\eta|=1$, with infinitesimal
generators $J_i=\tfrac{e_i}2+L_i$; harmonic and Fueter-regular functions decompose into the
irreducible summands $V_{2\ell+1}$ of this $SU(2)$-action, with quaternionic spherical harmonics
$\mathcal Y_\ell^m$ as angular factors \cite[\S\S5--6]{Baker2015}. The intrinsic (slice)
extension $R$ is the $\ell=0$ sector; raising and lowering operators $B_\pm$ move between
sectors, and $Q_\ell=(B_+)^\ell\circ Q_0$ (with $Q_0=-\tfrac14\Delta\circ R$ the Fueter map)
produces the regular representative of the $\ell$-th sector from a complex input
\cite[\S\S6--7]{Baker2015}. The key computation \cite[\S6]{Baker2016}: the image of the
exponential is
\begin{equation}\label{eq:Eell}
E_\ell(\lambda,z)=\frac{2}{\lambda^{2+\ell}}\,Q_\ell\bigl(e^{\lambda z}\bigr)
=i^\ell\bigl(j_\ell(\lambda y)+i\,j_{\ell+1}(\lambda y)\bigr)e^{\lambda x},
\end{equation}
an eigenfunction of the sector derivative with eigenvalue $\lambda$, by Rayleigh's formula.
So \emph{the spherical Bessel factor $j_\ell(\lambda y)$ is the radial profile with which a point
at height $y$ above the axis enters the $\ell$-th rotational sector at frequency $\lambda$}. The
manuscripts further develop the completeness of the $(\ell,m)$-decomposition, a Mellin/Laplace
inversion theory on each sector, and a fully quaternionic plane-wave expansion built from
\eqref{eq:Eell} with the addition theorem for the $\mathcal Y_\ell^m$
\cite[\S\S5--7]{Baker2016}; as a byproduct that manuscript derives a closed form for the
three-dimensional Clifford--Fourier kernel of
\cite{BrackxDeSchepperSommen2005,DeBieXu2011}, where closed forms had been known in even
dimensions only.

For present purposes the content is the following grounding of the definition
\eqref{eq:sector}. Apply the sector map $Q_\ell$ to the logarithmic derivative: for
$\xi'/\xi(z)=\sum_\rho(z-\rho)^{-1}$ (suitably paired), each zero at height $y$ above the axis
enters the $\ell$-th sector, at real frequency $\lambda$ of the sector Mellin/Laplace transform
\cite[Lemma~6.1, Thm.~6.2]{Baker2016}, with the radial profile $j_\ell(\lambda y)$ of
\eqref{eq:Eell}. The tags of \eqref{eq:sector} combine this radial profile with the unitary
axial character $e^{i\lambda a}$; rotating the radial profile as well produces the analytic
tags $i_\ell(\lambda y)e^{i\lambda a}$ of \S\ref{subsec:graded-anomaly} --- the two families
are the two real forms of one complexified family, matching at order $y^{2\ell}$. The sector
energies
\eqref{eq:sector} are thus the natural quadratic pairings
$\EE_\ell=\langle Q_\ell(\xi'/\xi),Q_\ell(\xi'/\xi)\rangle_w$ of the extension with itself ---
the $B_\pm$/Bessel machinery of \cite{Baker2015} put to quadratic use, as
Theorem~\ref{thm:linear-nogo} demands. We do not need the full strength of this dictionary
anywhere below: Lemma~\ref{lem:grading} and everything after it use \eqref{eq:sector} as a
definition, and this subsection explains why it is the canonical one.

\begin{remark}
The reader preferring to stay entirely within the slice-regular framework may take
\eqref{eq:sector} as a definition and skip \S\ref{sec:sectors}; nothing downstream depends on
it. Conversely, making the correspondence
$\{\text{sector energies}\}\leftrightarrow\{\text{spherical-harmonic content of the extended
zero set}\}$ into a theorem with a precise domain of validity (the natural route is the
plane-wave expansion of \cite[\S\S6--7]{Baker2016} applied to the Cauchy--Fueter kernel of the
extended $\xi'/\xi$) is worthwhile foundational work; it is listed with the open problems.
\end{remark}

%=====================================================================
\section{The explicit-formula bridge}\label{sec:bridge}

The unitary family of Section~\ref{sec:positivity} detects spheres from the zero side. To attack
RH one needs functionals computable from the \emph{prime} side. This section constructs the
analytic twin: a certified, prime-computable quadratic form whose growth rate reads off the
largest sphere radius.

\subsection{The certified explicit formula}\label{subsec:EF}

Rotated coordinates throughout: $\rho=\tfrac12+i\gamma$, the zero multiset $\{\gamma\}$ closed
under $\gamma\mapsto-\gamma$ (functional equation) and $\gamma\mapsto\bar\gamma$ (reality). For
$h$ even, holomorphic in $|\im r|\le\tfrac12+\varepsilon$, with sufficient decay, and
$g(u)=\tfrac1{2\pi}\int h(r)e^{-iur}\dd r$, the Guinand--Weil explicit formula
\cite{Guinand1948,Weil1952}, in the normalization of \cite[Thm.~5.12]{IwaniecKowalski2004}
specialized to $\zeta$ and rotated, reads
\begin{equation}\label{eq:EF}
\sum_\gamma h(\gamma)=h\bigl(\tfrac i2\bigr)+h\bigl(-\tfrac i2\bigr)
+\frac1{2\pi}\int_{-\infty}^\infty h(t)\Bigl[\re\,\psi\bigl(\tfrac14+\tfrac{it}2\bigr)
-\log\pi\Bigr]\dd t
-2\sum_{n\ge2}\frac{\Lambda(n)}{\sqrt n}\,g(\log n),
\end{equation}
with $\psi=\Gamma'/\Gamma$ and $\Lambda$ the von Mangoldt function. Because conventions differ
across sources, we certify \eqref{eq:EF} numerically as stated: on the Gaussian family below,
the zero side computed from $250$ zeros matches pole${}+{}$archimedean${}+{}$prime side to
$10^{-8}$--$10^{-10}$ at five $(c,\lambda)$ configurations (Appendix~\ref{app:numerics},
Table~\ref{tab:EF}).

\subsection{The linear building block and the quadratic form}

For a window center $c$, frequency $\lambda\ge0$ and width $\sigma$ (fixed at $\sigma=2$ in the
laboratory), take the even test function
\begin{equation}\label{eq:hfamily}
h_{c,\lambda}(r)=e^{i\lambda(r-c)}e^{-\frac{(r-c)^2}{2\sigma^2}}
+e^{-i\lambda(r+c)}e^{-\frac{(r+c)^2}{2\sigma^2}},
\qquad
Z(c,\lambda)=\sum_\gamma h_{c,\lambda}(\gamma).
\end{equation}
Its Fourier side is a Gaussian window at logarithmic scale $\lambda$:
\[
g(u)=\frac{\sigma}{\sqrt{2\pi}}\Bigl[e^{-iuc}e^{-\sigma^2(u-\lambda)^2/2}
+e^{iuc}e^{-\sigma^2(u+\lambda)^2/2}\Bigr],
\]
so by \eqref{eq:EF}, $Z(c,\lambda)$ is \emph{unconditionally computable from primes} with
$\log n\in[\lambda-O(1/\sigma),\lambda+O(1/\sigma)]$ plus explicit archimedean and pole terms:
the frequency $\lambda$ \emph{is} the logarithmic prime scale. The quadratic form is
$|Z(c,\lambda)|^2$ (pointwise, or integrated against a weight in $\lambda$).

\begin{lemma}[Diagonal anomaly]\label{lem:anomaly}
Fix $c>0$ and $\sigma$, and write $Y=\sup_\rho y_\rho$ for the zero multiset under
consideration, assumed to satisfy $y_\rho\le\tfrac12$ and $N(T)\ll T\log T$.
\begin{enumerate}
\item A real zero pair $\{\pm a\}$ contributes a bounded oscillatory amount to $Z(c,\lambda)$,
uniformly in $\lambda$.
\item A zero sphere $(a,y_0)$ --- the quadruple $\{\pm a\pm iy_0\}$ --- contributes to
$|Z(c,\lambda)|^2$ a diagonal block whose size oscillates, for large $\lambda$, between
\[
4(G-G')^2\,e^{2\lambda y_0}\quad\text{and}\quad4(G+G')^2\,e^{2\lambda y_0},
\qquad
G=e^{-\frac{(a-c)^2-y_0^2}{2\sigma^2}},\quad G'=e^{-\frac{(a+c)^2-y_0^2}{2\sigma^2}},
\]
up to $O(e^{\lambda y_0})$: exponential growth at rate exactly $2y_0$ in the $\limsup$ sense
(note $G>G'$ for $a,c>0$), weighted by the Gaussian distance of the sphere's center from the
window.
\item Unconditionally, $|Z(c,\lambda)|\le C_{c,\sigma}\,e^{\lambda Y}$, uniformly in $\lambda$;
in particular every product of two zero-terms in $|Z|^2$ is $O(e^{2\lambda Y})$, and the total
of all products in which at least one factor comes from a zero of height $\le Y-\varepsilon$ is
$O\bigl(e^{\lambda(2Y-\varepsilon)}\bigr)$.
\end{enumerate}
\end{lemma}

\begin{proof}
Write $h=h_{c,\lambda}$ as in \eqref{eq:hfamily}. (1) For real $\gamma$,
$|h(\gamma)|\le e^{-(\gamma-c)^2/2\sigma^2}+e^{-(\gamma+c)^2/2\sigma^2}\le2$, so a real pair
contributes at most $4$ to $|Z|$, uniformly in $\lambda$.

(2) Since $h$ is even, the quadruple contributes $2h(\gamma)+2h(\bar\gamma)$, $\gamma=a+iy_0$.
In $h(\bar\gamma)$ the first window term is
$e^{i\lambda(a-c)}e^{\lambda y_0}\cdot Ge^{i\theta}$ (with $\theta=(a-c)y_0/\sigma^2$ from the
complex Gaussian), and in $h(\gamma)$ the second window term is
$e^{-i\lambda(a+c)}e^{\lambda y_0}\cdot G'e^{i\theta'}$; all other terms have modulus
$e^{-\lambda y_0}\cdot O(1)$. Hence
\[
2h(\gamma)+2h(\bar\gamma)
=2e^{\lambda y_0}\Bigl(Ge^{i\lambda(a-c)+i\theta}+G'e^{-i\lambda(a+c)+i\theta'}\Bigr)
+O\bigl(e^{-\lambda y_0}\bigr),
\]
whose squared modulus oscillates between $4(G-G')^2e^{2\lambda y_0}$ and
$4(G+G')^2e^{2\lambda y_0}$, up to the stated error.

(3) Each zero $\gamma_k=a_k+iy_k$ satisfies $|h(\gamma_k)|\le
2e^{\lambda y_k}e^{y_k^2/2\sigma^2}e^{-(a_k-c)^2/2\sigma^2}$
(and the mirror term likewise). With $y_k\le\tfrac12$ the middle factor is bounded, and summing
the Gaussian factor against the zero-counting measure of density $O(\log(2+|a_k|))$ gives
$\sum_k|h(\gamma_k)|\le C_{c,\sigma}e^{\lambda Y}$, uniformly in $\lambda$. The refined bound
follows by using $e^{\lambda y_k}\le e^{\lambda(Y-\varepsilon)}$ on the low factor of each such
product.
\end{proof}

\begin{mainthm}[Rate criterion]\label{thm:rate}
Let the zero multiset satisfy the hypotheses of Lemma~\ref{lem:anomaly}, and set
$Y=\sup_\rho y_\rho$.
\begin{enumerate}
\item[(a)] \textup{(Unconditional.)} For every window $(c,\sigma)$,
$\limsup_{\lambda\to\infty}\lambda^{-1}\log|Z(c,\lambda)|\le Y$. In particular, if RH holds
then $Z(c,\lambda)$ is bounded in $\lambda$ for every window; contrapositively, exponential
growth of $Z$ in any single window disproves RH. $Z$ is unconditionally computable from primes
at logarithmic scale $\lambda$.
\item[(b)] \textup{(Exact rate. Hypothesis: either $Y=0$, or $Y>0$ and for some $\varepsilon>0$
only finitely many zeros satisfy $y_\rho>Y-\varepsilon$ --- a non-accumulation hypothesis; note
the $Y=0$ case must be listed separately, since when all zeros are real every zero satisfies
$y_\rho>-\varepsilon$ and the finiteness clause is vacuously false.)} For every fixed window,
\[
\limsup_{\lambda\to\infty}\frac{\log|Z(c,\lambda)|}{\lambda}
=\sup_\rho y_\rho=\Theta-\tfrac12,
\]
with $\Theta$ the supremum of real parts of the zeros of $\zeta$. Under this hypothesis:
\emph{RH holds if and only if $Z(c,\lambda)$ grows subexponentially in $\lambda$ for every
window}, and when the rate is positive it equals twice the supremal sphere radius, $2Y$, in the
$|Z|^2$ normalization.
\end{enumerate}
\end{mainthm}

\begin{proof}
(a) is part (3) of Lemma~\ref{lem:anomaly} (with $Y=0$, part (1) gives boundedness directly).

(b) If $Y=0$, part (a) gives $\limsup\le0$, and the mean-value argument below (applied to the
bounded, $\lambda$-independent-amplitude real-zero sum) shows $|Z|$ is bounded below along a
subsequence, so the $\limsup$ is exactly $0$. Now suppose $Y>0$ with the non-accumulation
hypothesis. Lower bound: fix $\varepsilon$ as given, let
$\gamma_1,\dots,\gamma_K$ (with their mirrors) be the finitely many zeros of height
$>Y-\varepsilon$, and split $Z=Z_{\mathrm{top}}+Z_{\mathrm{rest}}$ accordingly; part (3) of the
lemma gives $|Z_{\mathrm{rest}}|\le Ce^{\lambda(Y-\varepsilon)}$. By part (2), each
$\gamma_k=a_k+iy_k$ contributes to $Z_{\mathrm{top}}$ terms
$c_k^\pm e^{\lambda y_k}e^{\pm i\lambda(a_k\mp c)}$ with $\lambda$-independent amplitudes
$|c_k^\pm|>0$ plus $O(e^{-\lambda(Y-\varepsilon)})$, so
$e^{-\lambda Y}Z_{\mathrm{top}}$ is, up to $o(1)$ along $\lambda\to\infty$, a finite
exponential sum $\sum_k d_k e^{-\lambda(Y-y_k)}e^{i\lambda b_k}$ with distinct frequencies
$b_k$ after grouping (coincident frequencies group with coherent, non-cancelling coefficients:
the two growing contributions at frequency $a_k-c$ carry the same Gaussian phase). Let
$y_{\max}=\max_ky_k\in(Y-\varepsilon,Y]$. The Besicovitch mean value over windows
$[\Lambda,2\Lambda]$ of $|e^{-\lambda y_{\max}}Z_{\mathrm{top}}|^2$ tends to
$\sum_{k:y_k=y_{\max}}|d_k|^2>0$, so there is a sequence $\lambda_n\to\infty$ with
$|Z(c,\lambda_n)|\ge c'e^{\lambda_n y_{\max}}-Ce^{\lambda_n(Y-\varepsilon)}
\ge\tfrac{c'}2e^{\lambda_n(Y-\varepsilon)}$ eventually; hence
$\limsup\lambda^{-1}\log|Z|\ge Y-\varepsilon$. Shrinking $\varepsilon$ (the hypothesis for
$\varepsilon$ implies it for smaller $\varepsilon'$, with a possibly larger finite top set)
gives the equality with (a). The identification $Y=\Theta-\tfrac12$ is the
rotation $s=\tfrac12+iz$.
\end{proof}

\begin{remark}\label{rem:accumulation}
The hypothesis in (b) is automatic for the Davenport--Heilbronn function in
any fixed height range and for every finite or implanted configuration in the laboratory; for
$\zeta$ it holds under RH (the $Y=0$ branch), while under $\neg$RH the non-accumulation clause
would follow if the supremum $\Theta$ is attained, which is not
known. Removing it --- heights accumulating at the supremum, centers escaping to infinity ---
requires a multi-scale refinement of the mean-value argument that we have not carried out; we
flag this as a genuine, if technical, gap between (a) and the clean equivalence. Nothing in the
honest accounting of Remark~\ref{rem:honest} changes either way, and direction (a), which is the
falsifiable one, is unconditional.
\end{remark}

\textbf{Numerics} (Appendix~\ref{app:numerics}, Table~\ref{tab:rate}; Figure~\ref{fig:bridge}). With
$\sigma=2$, $\lambda\in[0.5,25]$, rate fitted on $\lambda>15$: $250$ real Riemann ordinates at
$c=20$ give rate $-0.0003$ ($|Z|^2=3.2$ at $\lambda=25$); the same data plus one implanted
sphere of radius $0.1$ at the window center give rate $0.105$
($|Z|^2=6.9\times10^2$) --- the fit is biased by the oscillatory background, since
$e^{2\lambda y_0}=e^5$ at the window edge barely clears it; the DH on-line zeros near $c=85.7$
alone give rate $-0.009$ ($|Z|^2=5.8$); and DH with its first sphere gives
\emph{rate $0.3085$ against the true $y_0=0.30852$} ($|Z|^2=2.1\times10^7$): the radius of a
genuine off-line zero recovered to four decimals from the growth of a single window.

\begin{remark}[Honest accounting]\label{rem:honest}
At the exponential-rate level, Theorem~\ref{thm:rate} is a Gaussian-mollified form of the
classical equivalence between zero-free regions and prime-counting error terms
($\psi(x)-x\ll x^{\Theta+\epsilon}$): detecting a sphere of radius $y_0$ through the growth of
$Z$ requires frequencies $\lambda\gtrsim1/y_0$, i.e.\ primes up to $e^{1/y_0}$. The wall is the
old wall; what the criterion adds is its geometry (window length $\leftrightarrow$ sphere
radius; growth rate $\leftrightarrow$ radius, exactly), a certified computational pipeline, and
a decomposition of the anomaly (below) over the same sectors as
Section~\ref{sec:positivity}. Genuinely new leverage, if any, lies below the exponential scale;
see Problems~\ref{prob:cone}--\ref{prob:matched}.
\end{remark}

\subsection{The graded decomposition of the anomaly}\label{subsec:graded-anomaly}

The sphere anomaly kernel $\cosh(2\lambda y_0)$ of Lemma~\ref{lem:anomaly} decomposes over the
quaternionic sectors with \emph{modified} Bessel tags: by the Gegenbauer addition theorem
(specialize \cite[10.60.2]{DLMF} at $u=v=z$, $\alpha=\pi$ and rotate by
\cite[10.47.12]{DLMF}; classically \cite[\S11.4]{Watson1944}),
\begin{equation}\label{eq:gegenbauer}
\frac{\sinh(2z)}{2z}=\sum_{\ell\ge0}(2\ell+1)\,i_\ell(z)^2 ,
\end{equation}
where $i_\ell$ is the modified spherical Bessel function of the first kind (verified to
$3\times10^{-14}$ at $z\in\{0.5,1,2,4\}$; Appendix~\ref{app:numerics}). The unitary family
($j_\ell$ tags, \S\ref{sec:positivity}) and the analytic family ($i_\ell$ tags) are Wick
rotations of each other --- $j_\ell(\lambda y)$ and $i_\ell(\lambda y)$ agree at leading order
$y^{2\ell}$ for small radius, and the analytic one is exponentially amplified once
$\lambda y_0\gtrsim1$. This is the same $j\leftrightarrow i$ rotation that runs through the
quaternionic extension itself: $\cos(qu)$ restricted to the sphere direction is a $\cosh$. Each
truncation of \eqref{eq:gegenbauer} splits the sphere anomaly into a positive series over
sectors, giving a family of inequalities interpolating between the Laguerre level ($\ell=1$) and
full Weil-type positivity; whether this graded family has a better cone than the classical one
is Problem~\ref{prob:cone}.

%=====================================================================
\section{The organizing statement: two involutions, and open problems}\label{sec:involutions}

\subsection{Two involutions}

The zero multiset of $\Xi$ (rotated coordinates, $\rho=\tfrac12+i\gamma$) carries two
involutions:
\[
\underbrace{\gamma\mapsto-\gamma}_{\text{holomorphic (functional equation, }
\rho\mapsto1-\rho)}
\qquad\text{and}\qquad
\underbrace{\gamma\mapsto\bar\gamma}_{\text{antiholomorphic (reality, }\rho\mapsto\bar\rho
\text{; the real structure }\sigma)} .
\]
The explicit formula \eqref{eq:EF} computes sums of functions \emph{analytic} in $\gamma$: it
natively sees the holomorphic involution (its test functions are taken even) and is structurally
blind to the antiholomorphic one. The sphere tags of Sections~\ref{sec:positivity} and
\ref{sec:bridge} are functions of the \emph{conjugate pair}: e.g.\ the analytic tag is exactly a
divided difference across the pair $(\gamma,\bar\gamma)$,
\[
i_0(\lambda y)\,e^{i\lambda a}
=\frac{e^{i\lambda\bar\gamma}-e^{i\lambda\gamma}}{-i\lambda(\gamma-\bar\gamma)},
\qquad\gamma=a+iy,
\]
and the unitary tag $j_0(\lambda y)e^{i\lambda a}$ is the same divided difference across the
rotated pair $(a+y,\,a-y)$ --- either way, the tags belong to the
class of \emph{conjugate-paired sums} $\sum_\gamma F(\gamma,\bar\gamma)$, strictly between
Weil-linear and fully quadratic. The sharpest scalar instance of the asymmetry:
\[
\sum_\rho\frac1{\rho(1-\rho)}=2+\gamma_{\mathrm{Euler}}-\log4\pi
\quad(\text{pairs by the functional equation; classical, see
\cite[Ch.~12]{Davenport2000}}),
\]
while $\sum_\rho1/(\rho\bar\rho)$ --- pairing by reality --- is not prime-computable by any
known identity; the termwise difference vanishes exactly on-line ($1-\rho=\bar\rho$), so the
defect between the two pairings is itself a sphere detector.

\begin{center}
\fbox{\parbox{0.92\textwidth}{\textbf{Organizing statement.} RH is the statement that the two
involutions coincide on the zero set: arithmetically, functional-equation pairing $=$
conjugation pairing; geometrically (Theorem~\ref{thm:dichotomy}), the zero divisor on the
incidence quadric is $\sigma$-invariant component-by-component. The program's central problem is
to import the antiholomorphic involution into explicit-formula arithmetic.}}
\end{center}

\subsection{Open problems}

Each problem carries its hardness-check annotation (\S\ref{sec:hardness}).

\begin{problem}[The graded cone]\label{prob:cone}
Does the $i_\ell$-graded decomposition \eqref{eq:gegenbauer} of the Weil-type quadratic form
define a positivity cone with better prime-side realizability, or better cross-term control,
than the classical Weil cone? \emph{(Must fail for DH --- it does structurally, DH having no
prime side; must degrade polynomially in $y_0$ --- it does, $y_0^{2\ell}$.)}
\end{problem}

\begin{problem}[Conjugate-paired explicit formulas]\label{prob:paired}
Find any unconditional identity, asymptotic, or one-sided bound for sums
$\sum_\gamma F(\gamma,\bar\gamma)$ beyond the trivial cases ($F$ depending on one variable, or
on $\gamma+\bar\gamma=2\re\gamma$ summed against known densities). Even partial results would
be new tools with reach beyond this program.
\end{problem}

\begin{problem}[Cross-term calculus]\label{prob:cross}
Quantify the off-diagonal of $|Z(c,\lambda)|^2$ (smoothed pair correlation; Montgomery's regime
\cite{Montgomery1973} is the on-line part) under RH$+$GUE and unconditionally; determine the
minimal sphere radius detectable at prime scale $e^\lambda$. This is the precise, quantitative
form of the wall of Remark~\ref{rem:honest} --- the gap any new idea must close.
\end{problem}

\begin{problem}[Coherent-mode detection]\label{prob:matched}
A sphere contributes a coherent $e^{\pm\lambda y_0}$-asymmetric mode to the pair-correlation
field, unlike GUE-type fluctuation. Formulate sphere-detection as matched filtering and
determine whether coherence buys detection below the $\lambda\sim1/y_0$ wall.
\emph{(Failure modes informative; must degrade polynomially near the axis.)}
\end{problem}

\begin{problem}[Equivariant/OCS route]\label{prob:ocs}
Compute local models of the lift and of the induced orthogonal complex structures at point-zeros
versus sphere-zeros (the technology of \cite{GSS2014,Altavilla2018,AltavillaSarfatti2019}
suffices for polynomial truncations), and determine whether Salamon--Viaclovsky-type
removability \cite{SalamonViaclovsky2009}, combined with order-$1$ growth and positivity of
$\Phi$, constrains $\sigma$-swapped components. \emph{(Must not be deformation-soft ---
Requirement 2.)}
\end{problem}

\begin{problem}[Spectral route]\label{prob:spectral}
The zero set of the intrinsic $\Xi$ is an axially symmetric set of exactly S-spectrum type
\cite{CSS2008,ColomboSabadiniStruppa2011}; quaternionic self-adjointness forces a real
S-spectrum. Realize $Z(\Xi)$ as the S-spectrum of a natural quaternionic operator built from the
explicit formula or from $\Phi$ --- a quaternionic Hilbert--P\'olya problem: same classical
difficulty, possibly better-shaped invariants (de Branges canonical systems are
quaternion-sized; the rotation about the real axis becomes a gauge action). Relatedly, make the
sector-energy/zero-set correspondence of \S\ref{sec:sectors} a theorem.
\end{problem}

%=====================================================================
\section{Formalization targets}\label{sec:formalization}

The program's foundations form a self-contained formalization target of independent value. As of
this writing (August 2026), mathlib \cite{mathlib} contains the quaternions as a normed
$\RR$-algebra with the expected analytic instances, but no slice analysis; the only formalization
of basic quaternionic (Fueter-regular) function theory we are aware of is in HOL Light
\cite{GabrielliMaggesi2017}. On the analytic side, mathlib has the completed Riemann zeta
function with its functional equation and Euler product (following the zeta/L-series
formalization program of Loeffler--Stoll \cite{LoefflerStoll2025}), and the Prime Number Theorem
has been proved in Lean by the PrimeNumberTheoremAnd project \cite{PNTand} (not yet merged into
mathlib at the time of writing).

Tier 1 (self-contained, valuable regardless of RH): slice regular basics, the representation
formula \eqref{eq:repformula}, the zero-structure theorem (Theorem~\ref{thm:zerostructure}), and
the reformulation (Theorem~\ref{thm:reformulation}) --- a new mathlib theory comparable in scope
to existing one-variable complex analysis developments. Tier 2: the incidence-quadric lift and
the equivariant statements of Section~\ref{sec:twistor}; the no-go theorems of
Section~\ref{sec:nogo}. Tier 3: the explicit formula \eqref{eq:EF} in a certified form ---
a natural continuation of \cite{LoefflerStoll2025,PNTand}. The numerics of
Appendix~\ref{app:numerics} (about $400$ lines of Python, reproducing in minutes) provide test
oracles for all formalized statements.

%=====================================================================
\appendix
\section{Numerical appendix}\label{app:numerics}

All computations reproduce from the ancillary files (\texttt{anc/xi\_lib.py},
\texttt{anc/run\_validation.py}, \texttt{anc/run\_experiments.py},
\texttt{anc/explicit\_formula.py}, \texttt{anc/run\_bridge.py}; Python with numpy, scipy,
mpmath, matplotlib; about $420$ lines total) in a few minutes on commodity hardware, plus a
one-time computation of $250$ zero ordinates via mpmath.
Figures~\ref{fig:lab} and \ref{fig:bridge} are generated by the included scripts.

\subsection{Validation of the extension and of DH (\texttt{run\_validation.py})}
$\Xi(x)=4\int_0^\infty\Phi(u)\cos(xu)\dd u$ ($6$ theta terms, trapezoid rule,
$du=5\times10^{-5}$ on $[0,3.2]$) matches $\xi(\tfrac12+ix)$ computed independently via mpmath
($\Gamma$, $\zeta$) to relative error $\le8\times10^{-10}$ at $x\in\{0,5,20,30\}$
($\xi(\tfrac12)=0.4971207781\ldots$; at $x\le20$ the agreement is $\sim10^{-14}$ or better),
and vanishes at the first zero ordinate
$x=14.134725141\ldots$ to $7\times10^{-18}$ (the quadrature floor). $\Xi_{DH}$ built from
Hurwitz zeta values is real and even in $z$ to relative $10^{-16}$ at $30$-digit working
precision (sampled); Newton iteration on $f_{DH}$ from crude seeds recovers the off-line zeros
\eqref{eq:DHspheres} with residuals $|f_{DH}|\le8\times10^{-30}$, matching \cite{Spira1994}.
The DH on-line ordinates in $[74,97]$, recovered by sign-change bisection of $\Xi_{DH}$ and
used in the rate experiment:
$74.325$, $75.909$, $77.226$, $78.717$, $80.328$, $81.925$, $83.109$, $87.647$, $89.439$,
$90.725$, $92.183$, $93.781$, $95.349$, $96.292$.

\subsection{Experiments 1--2: Laguerre detector (\texttt{run\_experiments.py})}
$L_1[\Xi](x)=\Xi'(x)^2-\Xi(x)\Xi''(x)$ with all three functions from the cosine transform:
the relative margin has minimum $0.173$ on $[0,52]$; beyond $x\approx55$ the float64 quadrature
floor ($|\Xi'|\sim10^{-18}$, $L_1\sim10^{-34}$) is reached and higher-precision quadrature would
be required for a verdict. For DH (rescaled $W=e^{\pi x/4}\Xi_{DH}$): $L_1<0$ exactly on
$[85.417,85.982]$ (predicted $[85.391,86.008]$) and $[114.025,114.302]$ (predicted
$[114.013,114.314]$); control $[78,82]$: $\min L_1\approx7.18\times10^2>0$.

\subsection{Experiment 3: sector energies (\texttt{run\_experiments.py})}
Weight $e^{-\lambda^2}$ on $\lambda\in[0,2]$; zero set $=20$ real Riemann ordinates plus an
implanted sphere at $a=28$, radius $y_0\in[10^{-3},10^{-0.5}]$: log--log slopes of $\EE_\ell$
versus $y_0$ (fitted over the six smallest radii, where the sphere term is unsaturated) are
$2.000/4.000/6.000$ for $\ell=1,2,3$; $\EE_1(\text{real ordinates only})=0$
exactly; per-sphere $\EE_1$ for the three DH spheres:
$4.360\times10^{-3}$, $1.062\times10^{-3}$, $2.593\times10^{-4}$ (successive ratios $4.11$,
$4.10$ vs.\ predicted $(y_0/y_0')^2=4.18$, $4.11$).

\subsection{Explicit-formula certification and rate experiment}
\begin{table}[h]
\centering
\caption{Certification of \eqref{eq:EF} on the family \eqref{eq:hfamily}, $\sigma=2$, $250$
zeros: zero side vs.\ pole${}+{}$archimedean${}+{}$prime side.}\label{tab:EF}
\begin{tabular}{lll}
\hline
$(c,\lambda)$ & zero side & $|\Delta|$ \\
\hline
$(20,0)$ & $1.86901176$ & $1.1\times10^{-10}$\\
$(20,1)$ & $0.96584078+1.42563531i$ & $1.2\times10^{-9}$\\
$(20,2)$ & $-0.85353747+1.53365468i$ & $6.4\times10^{-9}$\\
$(20,3)$ & $-1.80938739+0.21232192i$ & $1.8\times10^{-8}$\\
$(35,2)$ & $-0.38300486+0.14677138i$ & $4.8\times10^{-9}$\\
\hline
\end{tabular}
\end{table}

\begin{table}[h]
\centering
\caption{Rate experiment: $\sigma=2$, $\lambda\in[0.5,25]$, rate fitted to
$\log|Z|$ on $\lambda>15$.}\label{tab:rate}
\begin{tabular}{lccc}
\hline
zero set & $|Z|^2$ at $\lambda=25$ & fitted rate & true $\sup y_\rho$ \\
\hline
$\zeta$ zeros ($250$, real), $c=20$ & $3.2$ & $-0.0003$ & $0$\\
same $+$ implanted sphere $y_0=0.1$ at $a=c$ & $6.9\times10^2$ & $0.105$ & $0.1$\\
DH on-line ordinates only, $c=85.7$ & $5.8$ & $-0.009$ & $0$\\
DH with its first sphere, $c=85.7$ & $2.1\times10^7$ & $\mathbf{0.3085}$ & $\mathbf{0.30852}$\\
\hline
\end{tabular}
\end{table}

The implant fit is biased (here slightly high) because $e^{2\lambda y_0}=e^5$ at the window edge
barely clears the oscillatory background; the DH sphere, with $e^{2\lambda y_0}=e^{15.4}$, gives
the rate to four digits. The Gegenbauer identity \eqref{eq:gegenbauer} holds to
$\le3\times10^{-14}$ at $z\in\{0.5,1,2,4\}$ with $80$ terms.

\begin{figure}[p]
\centering
\includegraphics[width=\textwidth]{lab_results.pdf}
\caption{Left: $L_1[\Xi]>0$ on $[0,60]$ (semilog). Center: $L_1[\Xi_{DH}]$ dips negative
exactly inside the predicted windows $|x-a|<y_0$ at the first two DH spheres (shaded).
Right: sector energies $\EE_\ell\propto y_0^{2\ell}$, measured slopes $2.00/4.00/6.00$.}
\label{fig:lab}
\end{figure}

\begin{figure}[p]
\centering
\includegraphics[width=0.8\textwidth]{bridge_results.pdf}
\caption{The rate criterion (Theorem~\ref{thm:rate}) on real data: $|Z(c,\lambda)|^2$ versus
$\lambda$ for $250$ real Riemann ordinates ($c=20$), the same plus an implanted sphere of radius
$0.1$, the DH on-line ordinates near $c=85.7$, and DH with its first sphere ($y_0=0.30852$;
fitted rate $0.3085$).}
\label{fig:bridge}
\end{figure}

%=====================================================================
\begin{thebibliography}{99}

\bibitem{Altavilla2018}
A.~Altavilla,
\emph{Twistor interpretation of slice regular functions},
J. Geom. Phys. \textbf{123} (2018), 184--208.

\bibitem{AltavillaSarfatti2019}
A.~Altavilla and G.~Sarfatti,
\emph{Slice-polynomial functions and twistor geometry of ruled surfaces in $\mathbb{CP}^3$},
Math. Z. \textbf{291} (2019), no. 3--4, 1059--1092.

\bibitem{Baker2015}
E.~B. Baker III,
\emph{Quaternionic analysis with applications to the Riemann zeta function},
manuscript (2015); archived in \cite{ZenodoArchive}.

\bibitem{Baker2016}
E.~B. Baker III,
\emph{The quaternion-regular exponential and the three dimensional Clifford--Fourier kernel},
manuscript (2016); archived in \cite{ZenodoArchive}.

\bibitem{Baker2017}
E.~B. Baker III,
\emph{Twistor holomorphicity},
manuscript (2017); archived in \cite{ZenodoArchive}.

\bibitem{Baker2022arXiv}
E.~B. Baker III,
\emph{Causal and self-dual morphisms in four complex dimensions},
arXiv:2203.07952 (2022).

\bibitem{Baker2023arXiv}
E.~B. Baker III,
\emph{Generalized conformal maps as classical symmetries of Yang--Mills fields},
arXiv:2301.00029 (2023).

\bibitem{Baker2023}
E.~B. Baker III,
\emph{Frustrated conformal transformations and holomorphic maps on ambitwistor space},
manuscript (2023); archived in \cite{ZenodoArchive}.

\bibitem{ZenodoArchive}
E.~B. Baker III,
\emph{The zero-sphere program: archival manuscripts, working memos, and code (2015--2026)},
Zenodo, DOI: [to appear].

\bibitem{BalanzarioSanchez2007}
E.~P. Balanzario and J.~S\'anchez-Ortiz,
\emph{Zeros of the Davenport--Heilbronn counterexample},
Math. Comp. \textbf{76} (2007), no. 260, 2045--2049.

\bibitem{BoryReyesShapiro2010}
J.~Bory Reyes and M.~Shapiro,
\emph{Clifford analysis versus its quaternionic counterparts},
Math. Methods Appl. Sci. \textbf{33} (2010), no. 9, 1089--1101.

\bibitem{BrackxDeSchepperSommen2005}
F.~Brackx, N.~De Schepper, and F.~Sommen,
\emph{The Clifford--Fourier transform},
J. Fourier Anal. Appl. \textbf{11} (2005), no. 6, 669--681.

\bibitem{CGSS2009}
F.~Colombo, G.~Gentili, I.~Sabadini, and D.~C. Struppa,
\emph{Extension results for slice regular functions of a quaternionic variable},
Adv. Math. \textbf{222} (2009), no. 5, 1793--1808.

\bibitem{CSS2008}
F.~Colombo, I.~Sabadini, and D.~C. Struppa,
\emph{A new functional calculus for noncommuting operators},
J. Funct. Anal. \textbf{254} (2008), no. 8, 2255--2274.

\bibitem{ColomboSabadiniStruppa2011}
F.~Colombo, I.~Sabadini, and D.~C. Struppa,
\emph{Noncommutative Functional Calculus: Theory and Applications of Slice Hyperholomorphic
Functions},
Progress in Mathematics \textbf{289}, Birkh\"auser, Basel, 2011.

\bibitem{CsordasNorfolkVarga1986}
G.~Csordas, T.~S. Norfolk, and R.~S. Varga,
\emph{The Riemann hypothesis and the Tur\'an inequalities},
Trans. Amer. Math. Soc. \textbf{296} (1986), no. 2, 521--541.

\bibitem{CsordasVarga1990}
G.~Csordas and R.~S. Varga,
\emph{Necessary and sufficient conditions and the Riemann hypothesis},
Adv. in Appl. Math. \textbf{11} (1990), no. 3, 328--357.

\bibitem{CsordasVishnyakova2013}
G.~Csordas and A.~Vishnyakova,
\emph{The generalized Laguerre inequalities and functions in the Laguerre--P\'olya class},
Cent. Eur. J. Math. \textbf{11} (2013), no. 9, 1643--1650.

\bibitem{Davenport2000}
H.~Davenport,
\emph{Multiplicative Number Theory}, 3rd ed., revised by H.~L. Montgomery,
Graduate Texts in Mathematics \textbf{74}, Springer, New York, 2000.

\bibitem{DavenportHeilbronn1936a}
H.~Davenport and H.~Heilbronn,
\emph{On the zeros of certain Dirichlet series (I)},
J. London Math. Soc. \textbf{11} (1936), 181--185.

\bibitem{DavenportHeilbronn1936b}
H.~Davenport and H.~Heilbronn,
\emph{On the zeros of certain Dirichlet series (II)},
J. London Math. Soc. \textbf{11} (1936), 307--312.

\bibitem{DeBieXu2011}
H.~De Bie and Y.~Xu,
\emph{On the Clifford--Fourier transform},
Int. Math. Res. Not. IMRN \textbf{2011}, no. 22, 5123--5163.

\bibitem{DLMF}
\emph{NIST Digital Library of Mathematical Functions}, \url{https://dlmf.nist.gov/},
Release 1.2.x; F.~W.~J. Olver et al., eds.

\bibitem{GabrielliMaggesi2017}
A.~Gabrielli and M.~Maggesi,
\emph{Formalizing basic quaternionic analysis},
in: Interactive Theorem Proving (ITP 2017), Lecture Notes in Comput. Sci. \textbf{10499},
Springer, 2017, pp. 225--240.

\bibitem{GentiliStoppato2008}
G.~Gentili and C.~Stoppato,
\emph{Zeros of regular functions and polynomials of a quaternionic variable},
Michigan Math. J. \textbf{56} (2008), no. 3, 655--667.

\bibitem{GSSbook}
G.~Gentili, C.~Stoppato, and D.~C. Struppa,
\emph{Regular Functions of a Quaternionic Variable}, 2nd ed.,
Springer Monographs in Mathematics, Springer, Cham, 2022.

\bibitem{GentiliStruppa2007}
G.~Gentili and D.~C. Struppa,
\emph{A new theory of regular functions of a quaternionic variable},
Adv. Math. \textbf{216} (2007), no. 1, 279--301.

\bibitem{GSS2014}
G.~Gentili, S.~Salamon, and C.~Stoppato,
\emph{Twistor transforms of quaternionic functions and orthogonal complex structures},
J. Eur. Math. Soc. (JEMS) \textbf{16} (2014), no. 11, 2323--2353.

\bibitem{GhiloniPerotti2011}
R.~Ghiloni and A.~Perotti,
\emph{Slice regular functions on real alternative algebras},
Adv. Math. \textbf{226} (2011), no. 2, 1662--1691.

\bibitem{GORZ2019}
M.~Griffin, K.~Ono, L.~Rolen, and D.~Zagier,
\emph{Jensen polynomials for the Riemann zeta function and other sequences},
Proc. Natl. Acad. Sci. USA \textbf{116} (2019), no. 23, 11103--11110.

\bibitem{GriffithsHarris1978}
P.~Griffiths and J.~Harris,
\emph{Principles of Algebraic Geometry},
Wiley, New York, 1978.

\bibitem{GriffithsHarris1979}
P.~Griffiths and J.~Harris,
\emph{Algebraic geometry and local differential geometry},
Ann. Sci. \'Ecole Norm. Sup. (4) \textbf{12} (1979), 355--432.

\bibitem{Guinand1948}
A.~P. Guinand,
\emph{A summation formula in the theory of prime numbers},
Proc. London Math. Soc. (2) \textbf{50} (1948), 107--119.

\bibitem{IwaniecKowalski2004}
H.~Iwaniec and E.~Kowalski,
\emph{Analytic Number Theory},
American Mathematical Society Colloquium Publications \textbf{53},
Amer. Math. Soc., Providence, RI, 2004.

\bibitem{LiuMorenoShi2026}
C.~Liu, G.~Moreno, and H.~Shi,
\emph{Classifying slice-regular polynomials via group actions on the twistor space},
arXiv:2605.22788 (2026).

\bibitem{LoefflerStoll2025}
D.~Loeffler and M.~Stoll,
\emph{Formalizing zeta and L-functions in Lean},
Ann. Formalized Math. \textbf{1} (2025); arXiv:2503.00959.

\bibitem{mathlib}
The mathlib Community,
\emph{The Lean mathematical library},
\url{https://github.com/leanprover-community/mathlib4}.

\bibitem{MasonWoodhouse1996}
L.~J. Mason and N.~M.~J. Woodhouse,
\emph{Integrability, Self-Duality, and Twistor Theory},
London Mathematical Society Monographs, Clarendon Press, Oxford, 1996.

\bibitem{Montgomery1973}
H.~L. Montgomery,
\emph{The pair correlation of zeros of the zeta function},
in: Analytic Number Theory (Proc. Sympos. Pure Math. XXIV, St. Louis, 1972),
Amer. Math. Soc., Providence, RI, 1973, pp. 181--193.

\bibitem{PNTand}
A.~Kontorovich, T.~Tao, et al.,
\emph{PrimeNumberTheoremAnd},
Lean formalization project,
\url{https://github.com/AlexKontorovich/PrimeNumberTheoremAnd}.

\bibitem{Polya1926}
G.~P\'olya,
\emph{Bemerkung \"uber die Integraldarstellung der Riemannschen $\xi$-Funktion},
Acta Math. \textbf{48} (1926), no. 3--4, 305--317.

\bibitem{Polya1927}
G.~P\'olya,
\emph{\"Uber die algebraisch-funktionentheoretischen Untersuchungen von J.~L.~W.~V. Jensen},
Kgl. Danske Vid. Sel. Math.-Fys. Medd. \textbf{7} (1927), 3--33.

\bibitem{RodgersTao2020}
B.~Rodgers and T.~Tao,
\emph{The de Bruijn--Newman constant is non-negative},
Forum Math. Pi \textbf{8} (2020), e6, 62 pp.

\bibitem{SalamonViaclovsky2009}
S.~Salamon and J.~Viaclovsky,
\emph{Orthogonal complex structures on domains in $\mathbb{R}^4$},
Math. Ann. \textbf{343} (2009), no. 4, 853--899.

\bibitem{Spira1994}
R.~Spira,
\emph{Some zeros of the Titchmarsh counterexample},
Math. Comp. \textbf{63} (1994), no. 208, 747--748.

\bibitem{Titchmarsh1986}
E.~C. Titchmarsh,
\emph{The Theory of the Riemann Zeta-Function}, 2nd ed.,
edited and with a preface by D.~R. Heath-Brown,
Oxford University Press, 1986.

\bibitem{Watson1944}
G.~N. Watson,
\emph{A Treatise on the Theory of Bessel Functions}, 2nd ed.,
Cambridge University Press, 1944.

\bibitem{Weil1952}
A.~Weil,
\emph{Sur les ``formules explicites'' de la th\'eorie des nombres premiers},
Comm. S\'em. Math. Univ. Lund (1952), Tome Suppl\'ementaire, 252--265.

\bibitem{Zak1993}
F.~L. Zak,
\emph{Tangents and Secants of Algebraic Varieties},
Translations of Mathematical Monographs \textbf{127},
Amer. Math. Soc., Providence, RI, 1993.

\end{thebibliography}

\end{document}
